\documentclass{report}
\usepackage{polyglossia}
\usepackage[backend=biber,
style=verbose,
sorting=ynt]{biblatex}
\addbibresource{bib.bib}
\setdefaultlanguage{french}
\usepackage{hyperref}
\usepackage{graphicx}
\usepackage{verbatim}
\usepackage[a4paper,
inner=3cm,T
outer=3cm,
top=3cm,
bottom=3cm,
]{geometry}
\usepackage{xcolor}
\newcommand{\og}[0]{\guillemetleft}
\newcommand{\fg}[0]{\guillemetright~}
\hypersetup{
    colorlinks,
    linkcolor={red!50!black},
    citecolor={blue!50!black},
    urlcolor={blue!80!black}
}

\title{Piratage culturel et communautés patrimoniales\\\vspace{0.5cm}Master Archive - Technologies numériques appliquées à l'histoire\\École nationale des chartes}
\author{Jules Musquin\\sous la direction d'Emmanuelle Bermès}

\date{2024-2025\\ \vspace{1cm}\begin{quote}\textit{[...] Lent dans mon ombre, j'explore la pénombre profonde, d'une canne indécise, moi qui m'imaginais le Paradis sous l'espèce d'une bibliothèque.}\\\vspace{0.2cm}Jorge Luis Borges, \textit{Poème des dons}\end{quote}}
\begin{document}

\maketitle

\chapter*{}

\textit{Merci à Mme Emmanuelle Bermès, qui m'a fait confiance pour travailler sur ce sujet qui me passionne, et qui m'a dirigée et accompagnée tout au long de cette année.
}\\

\textit{Merci à l'Association des archivistes français pour avoir aidé à financer cette année d'étude 2024-2025.}\\

\textit{Merci à tous mes enquêtés, dont les personnalités et expériences fondent une grande partie de ce travail. Plus largement, je souhaite adresser des remerciements à l’ensemble de ma promotion pour les échanges enrichissants auxquels nos sujets de mémoire respectifs ont donné lieu.}\\

\textit{Merci à ceux qui font vivre et maintiennent les bibliothèques pirates sur lesquelles je passe tant de temps.}

\tableofcontents


\chapter{Introduction}

% Accroche & introduction du sujet

\begin{quote}
    \textit{Mais dès l'instant où le critère de l'authenticité s'avère inapplicable à la production artistique, la fonction sociale entière de l'art s'en trouve bouleversée. Au lieu de se fonder sur le rituel, elle se fonde sur une autre praxis : en l'occurrence, sur la politique.}
    
    Walter Benjamin, \textit{L'\oe uvre d'art à l'époque de sa reproductibilité technique}, 1939.
\end{quote}

L'acte de copie d'un objet culturel est la première étape à son partage, sa conservation et sa réinterprétation. Si une langue a besoin de locuteurs pour pouvoir évoluer, alors un objet culturel qui n'est pas copié ne peut pas être nommé comme un \guillemetleft objet culturel \guillemetright ~, tout comme le latin ne peut pas être considéré comme une \guillemetleft langue vivante \guillemetright. Un premier exemple de l'importance de la copie et par extension du piratage dans notre patrimoine est \textit{Le Voyage dans la Lune} réalisé en 1902 par Georges Méliès. La plupart des copies officielles de ce film ont disparu, ou ont été détruites par Méliès lui-même alors qu'il ne pouvait plus les entreposer à cause de la fermeture de son studio. Les chercheurs et restaurateurs ont pu reconstituer et préserver son œuvre grâce à de très nombreuses copies qui étaient soit piratées, soit confisquées \footnote{Marie-France Briselance et Jean-Claude Morin, \textit{Grammaire du cinéma}, Paris, Nouveau Monde, coll. « Cinéma », 2010, 588 p. (ISBN 978-2-84736-458-3), p. 111.}.

Un deuxième exemple, peut-être le plus connu, est celui du film \textit{Nosferatu le vampire} réalisé en 1922 par Friedrich Wilhelm Murnau. Florence Stoker, veuve de Bram Stoker, l'auteur du roman \textit{Dracula}, intente à Murnau un procès pour plagiat que ce dernier perd. Toutes les copies ainsi que les négatifs sont détruits par deux fois en 1925 et 1928. Rapidement, d'anciennes copies cachées refont surface - et le film survit grâce à des projections confidentielles jusqu'en 1981 où plusieurs copies partielles permettent de le restaurer au plus près de l'archétype du travail de Murnau \footnote{Jean Pierre Piton, \textit{L'éternel retour de Dracula}, L'Écran Fantastique no 130 janvier 1993 p. 21.}.

La numérisation de nos outils professionnels, de nos sources de divertissement et de nos modes de communication, généralise la capacité de copie des objets culturels qui peuplent nos vies. Cette propriété préexistait déjà avec, par exemple, la K7 VHS à bande magnétique qui permet d'enregistrer un signal analogique. Cependant, dans un objet numérique, la propriété de copie est parfaitement intégrée au fonctionnement du fichier et de son système d'accès par défaut, l'ordinateur, produisant un dilemme numérique entre la liberté du partage d'une information réplicable gratuitement et la propriété intellectuelle littéraire et artistique \footnote{\cite{national_research_council_us_digital_2010}}.

La nouvelle facilité de copie et d'échange de l'information grâce à la généralisation de l'informatique bon marché et de l'accès au web domestique a rapproché des individus qui se sont organisés en communautés culturelles,  partageant des références culturelles et les fichiers qui les documentent, devenant ainsi des communautés patrimoniales.

Les communautés patrimoniales \footnote{~Convention de Faro sur la valeur du patrimoine culturel pour la société, 2005} qui nous intéressent dans ce travail sont celles qui décident d'embrasser la copie et le partage de l'information dans sa forme la plus radicale, en privilégiant les \guillemetleft communs numériques \guillemetright. Elles peuvent choisir de produire et d'entretenir leur propre contenu à la manière d'un Wikipédia, devenant des \textit{communautés libres}, ou elles peuvent choisir de faire fi des règles et des lois en devenant des \textit{communautés patrimoniales pirates}\footnote{\cite{bermes_lecran_2024}}. Si les communautés libres essaient de se différencier le plus possible des communautés pirates, les individus les composant évoluent entre les deux, puisant dans les ressources des communautés libres pour ensuite compléter leurs besoins avec l'aide du piratage.

Nous étudierons ici les liens, les émotions, et les mouvements de va-et-vient entre les individus qui composent les communautés pirates, les objets numériques piratés, et les lieux virtuels où se construisent ces communautés. Notre premier but est d'étendre le concept de communautés patrimoniales aux pirates culturels, puis de comprendre l’appréhension du public du \guillemetleft dilemme numérique de la propriété intellectuelle \guillemetright, et enfin d'étudier les pratiques émergentes de partage et de conservation numériques qui mobilisent ces communautés patrimoniales pirates.

\section{Langues et jambes de bois : définir la piraterie culturelle}

Les travaux sur la véritable piraterie historique ont souvent une introduction permettant de séparer une partie des représentations du sujet et la réalité des faits. Je pense que pour la piraterie culturelle, un travail similaire est le bienvenu. Les travaux sur le piratage se répartissent sur deux axes bibliographiques principaux.\\

\subsection*{La piraterie culturel entre les enjeux scientifiques et la communication}

Les travaux sur le piratage se répartissent sur deux axes bibliographiques principaux.\\

Le premier se concentre sur les industries de la propriété intellectuelle, essentiellement dans le domaine du marché pharmaceutique international. Pour ce qui est de la piraterie de biens culturels, le sujet le plus étudié est la place du piratage dans les économies des pays en développement, à la fois ateliers du monde du capitalisme planétaire, mais aussi envisagé comme une manne de consommateurs potentiels des divertissements du nord économique. Des revues de science de l'information, de communication, de management et d'économie étudient l'organisation et les chemins de la distribution du contenu culturel des pays du Sud économique comme par exemple à Malégaon en Inde \footnote{Lawrence Liang,\textit{ Piratage, créativité et infrastructure : repenser l’accès à la culture}, Tracés. Revue de Sciences humaines, 2014}, ou au Cameroun \footnote{Mattelart, T.\textit{ Chapitre 1. Piratages audiovisuels et réseaux de la mondialisation par le bas}. Dans \og Piratages audiovisuels : Les voies souterraines de la mondialisation culturelle \fg, p. 27-52, éditions De Boeck Supérieur, 2011}.

Le deuxième se concentre sur la piraterie culturelle numérique comme un marché, où le monde académique tente de saisir à la fois les lieux où s'organisent les échanges, et de comprendre comment s'organisent ceux qui font vivre ces lieux. Le mot de \guillemetleft  Pirate \guillemetright~, ainsi que l'iconographie de la piraterie historique, est également un point de rendez-vous, ou un repoussoir - pour les communautés qui investissent le web comme un espace d'échange de biens culturels. L'exemple le plus partagé est le nom et l'iconographie du site \textit{The Pirate Bay}, dont l'icône est un trois-mâts noir, dont la voile centrale montre une variation du \textit{Jolly Roger} remplaçant le crâne par une cassette en inscrivant en dessous dans une police gothique le nom du site de partage de torrents le plus connu des internets. L'utilisation d'autres termes comme \textit{partage illégal du contenu culturel, reproduction interdite ou non autorisée, partage d'informations avec son voisin} est une alternative qui offre une représentation plus mesurée et plus proche de la réalité des faits \footnote{Richard Stalman,\textit{Quelques expressions trompeuses qu'il vaut mieux éviter d'employer}, dans \guillemetleft Libre enfants du savoir numérique \guillemetright, p.359-364 }. Au-delà de cette analyse légaliste et quantitative, la sphère académique des sciences humaines et sociales s'est intéressée à la popularité de plateformes telles que Napster \footnote{Moulier-Boutang, \textit{Droits de propriété intellectuelle, terra nullius et capitalisme cognitif}. Multitudes,n°41(2), p66-72, 2010} ou The Pirate Bay . Ces recherches universitaires ont pour but de comprendre la psychologie des individus en examinant les discours et valeurs associés au piratage.

\subsection*{La communauté patrimoniale pirate}

L'objet numérique et son partage mobilisent des concepts et des catégories juridiques diverses qu'il est important de définir. Cet objet numérique s'incorpore dans le concept de \textit{patrimoine documentaire immatériel}, promu et adopté par la Convention de l'UNESCO pour que ses États membres puissent intégrer au droit des définitions et des outils politiques de promotion de la diversité culturelle, y compris de la culture numérique et de la culture du web. Si un patrimoine est immatériel, c'est alors qu'il est constitué d'objets qui ont un sens particulier pour un groupe de personnes, le patrimoine immatériel étant donc la relation entre ces \textit{objets} et les groupes de personnes qui leur donnent un sens et une fonction. Quand un objet est la cause principale de lien dans un groupe (et non plus une conséquence de ce lien), nous pouvons alors parler de l'apparition d'une \textit{communauté patrimoniale}. Ce terme de \textit{communauté patrimoniale} entre dans le vocabulaire institutionnel en 2005 avec la signature de la convention européenne de Faro sur la valeur du patrimoine culturel pour la société. Cette convention inscrit l'action européenne dans le domaine du patrimoine avec un nouveau paradigme, considérant autant les objets et les lieux que les usages, les personnes et les interactions qui en résultent.

Les émotions négatives sont une composante puissante dans la définition de ce qu'est le patrimoine (peur de voir disparaître, de voir changer, de voir trahir un objet), les communautés se soudent avec des émotions comme la colère, la sédition, la rébellion face à l'injustice \footnote{\cite{fabre_emotions_2013}}, en invoquant des référentiels politiques ou institutionnels.  Les communautés patrimoniales peuvent entrer dans des dynamiques d'actions de sauvegarde de leur patrimoine, parfois avec la volonté d'aider les acteurs institutionnels ou entrepreneuriaux (une restauration d'un bâtiment par exemple), ou en entrant dans un rapport de force avec ces mêmes acteurs (la mobilisation des mouvements écologistes face à des travaux d'infrastructures comme les bassines de Sainte-Soline en est un bon exemple). 

Si le partage non autorisé de biens culturels n'a rien de nouveau, le tournant numérique de nos sociétés a engendré deux conséquences. La première est la massification des échanges, des capacités de stockage et de copie des individus,  la deuxième est la capacité nouvelle des individus à s'affranchir du temps et de l'espace pour entrer en relation, donnant ainsi naissance à des communautés patrimoniales plus nombreuses et beaucoup plus développées dès la démocratisation de l'accès au web. Quand une communauté patrimoniale se forme autour d'objets ou de pratiques \textit{illégales} nous pouvons alors utiliser le terme de communautés patrimoniales \textit{pirates}. Notre travail a pour objectif de mieux comprendre les relations entre l'objet, l'individu et la communauté dans ce cadre de sortie des référentiels institutionnels et entrepreneuriaux - le but des pirates n'étant pas de trouver une légitimité aux yeux de l'État ou des entreprises, mais bien de s'en affranchir.

\section{Revue bibliographique}

Étudier des individus qui piratent implique de se plonger dans leurs histoires personnelles, et les origines de leur montée en compétence progressive - passant de néophyte à amateur, puis pour certains à des professionnels du piratage, ce qu'Anthony Galluzzo, professeur en science de gestion à l'université de Saint-Étienne, nomme une \textit{carrière pirate}\footnote{\cite{galluzzo_longevite_2021}}. Cette acquisition de compétences techniques est parallèle à une éducation, un apprentissage social et politique. C'est en combinant les deux que les pirates génèrent une sociabilité particulière s'appuyant autant sur des éléments techniques (comment accéder à une ressource, qui possède un tutoriel pour tel outil logiciel, qui possède une copie d'un objet numérique rare que je cherche) que sur des émotions patrimoniales (ce site va disparaître ; comment aider ma communauté à le préserver). 

Ce lien a été l'objet d'étude de la littérature anglo-saxonne des \textit{fan studies} \footnote{\cite{taylor_save_nodate}} \footnote{:\cite{de_kosnik_rogue_2016}}. Ces études se concentrent sur les cultures du web (des fanfictions, le jeu vidéo) et la manière dont les personnes qui vivent ces cultures comprennent et organisent la mémoire de ce qui a été produit au sein de leur communauté. Nous étudions une culture \textit{nativement numérique}, mais certaines différences sur la situation des communautés patrimoniales pirates ne permettent pas de plaquer les analyses des \textit{fantstudies} sur ce sujet.

L'un des paramètres ayant le plus d'influence sur l'organisation des communautés patrimoniales pirates et de leurs membres est \textit{l'organisation anarchique de guérilla}. Comme pour Wikipédia, la plupart des dépôts communautaires pirates n'obéissent pas à une structure verticale et rigide. Un noyau dur de membres reconnus pour leur implication et leur ancienneté dans le projet permet de maintenir l'infrastructure de base du dépôt communautaire, ce à quoi s'ajoutent des membres moins investis qui offrent leur aide et leurs conseils de manière ponctuelle. Un troisième cercle est la communauté large autour du dépôt, des usagers réguliers offrant un soutien financier, qui propagent la bonne adresse à l'aide du bouche-à-oreille sur les forums, ou simplement qui témoignent de l'importance du dépôt pirate dans leur vie de tous les jours. 

Le terme de guérilla fait référence à la manière dont ces dépôts, leurs membres et leurs usagers s'organisent pour mettre en place un espace à la fois efficace dans sa mission (transmettre et archiver des données) sans se faire repérer par les éditeurs ou les autorités de protection des droits d'auteur \footnote{Martin Paule Eve, Lessons from the Library: Extreme Minimalist Scaling at Pirate Ebook Platforms, dans \textit{Digital Humanities Quarterly} Vol 16 n°2, 2022}. Pour cela, les dépôts pirates doivent à la fois avoir une visibilité minimale (ne pas être vus), sans que le trafic internet puisse être retracé à l'hébergeur du dépôt (ne pas être visé), avec une résilience aux attaques et aux fermetures d'hébergeur (ne pas être touchés).

\section{Nos sources}

Une partie de notre étude se base sur l'étude des archives et des forums de communautés patrimoniales pirates. Une partie a fait l'objet d'un archivage, parfois d'une patrimonialisation et d'une curation, c'est par exemple le cas des archives du web de la BnF ou des collections de The Internet Archive. Une autre partie de nos sources a été directement archivée en utilisant l'outil d'instantané (snapshot) Zotero.
\begin{itemize}
    \item Nous avons effectué un travail de recherche sur les archives du web de la BnF, en particulier les archives du site Skyblog. Grâce à Marina Hervieu, chargé de projet de recherche à la BnF, nous avons pu explorer les skyblogs ayant dans leurs pseudonymes des mots laissant penser à une pratique régulière du piratage \footnote{pirate, piratage, téléchargement, télécharger, torrenting, torrent, fansub, mininova, des mots ou des noms de service en lien avec le piraterie durant les années de service de Skyblog}. Dans les listes de blog minées par Mme Hervieu, nous avons ensuite pu appliquer un tri en fonction de la popularité et du nombre de posts du blog. En utilisant cette méthode, nous avons très rapidement atteint une saturation dans nos d'exemples. 
    
    \item Une grande proportion de nos sources provient de l'observation des archives et de l'actualité des forums des communautés patrimoniales pirates, notamment sur Reddit, Twitter/X et Discord. Les deux sites les plus propices à l'étude se sont révélés être Discord et Reddit, du fait que la plateforme Twitter/X n'a introduit la possibilité de former des « communautés » que très récemment. 
    
    \subitem Discord est un système de messagerie communautaire où nous avons suivi la communauté qui se nomme \textit{The Eye}. C'est un site communautaire issu de la communauté \textit{Archive Team}, ayant pour but de préserver et de partager des pans de la culture numérique. 

    \begin{quote}
        \textit{Living without the knowledge of our past history, origin and culture is like a tree without roots. The Internet is a worldwide platform for sharing information. It is a community of common interests.
        Our mission at The Eye is to preserve pieces of digital history. We are digital librarians.}
        
        Message de présentation du serveur discord de \textit{The Eye}
    \end{quote}

    \subitem le site Reddit est également une source très instructive pour étudier les communautés patrimoniales pirates. Il est composé de \textit{subreddits}, des forums communautaires sur des sujets spécifiques. Il existe plusieurs subreddit sur les sujets du piratage de bien culturel et de la sauvegarde du patrimoine numérique. Bien que l'intégralité d'un subreddit ne puisse être classifiée comme une communauté patrimoniale, certaines communautés investissent activement des subreddits spécifiques. Nous pouvons illustrer cela par \textit{r/datahoarder} (des individus collectant de vastes quantités de données numériques) qui sert de point central pour les discussions et la mobilisation des amateurs d'archives numériques et de patrimoine dématérialisé, ou r/Piracy qui sert de point de ralliement pour toutes les discussions concernant le piratage.

    \subitem Twitter/X est une source plus difficilement exploitable pour notre travail car il s'agit d'une application de microblog où les gens peuvent se suivre de manière individuelle sans avoir l'option d’agréger les utilisateurs en communautés\footnote{Il existe depuis la reprise de l'application par Elon Musk un onglet communauté, mais la perte de popularité de Twitter couplé à l'ajout de cette fonction relativement peu utilisée fait qu'il reste une application de microblog individualisée. }. Les sources que nous agrégeons de Twitter ne sont pas représentatives d'une communauté patrimoniale, cependant certains sujets transversaux peuvent mobiliser un grand groupe d'utilisateurs, créant ainsi un effet de mode et d’entraînement permettant d'étudier l'importance de certains évènements sur la création de communautés patrimoniales pirates. L'exemple le plus récent dont je peux témoigner est la mort du réalisateur américain David Lynch. Les \textit{twittos}\footnote{Twittos : surnom vulgaire donné à un utilisateur intensif du réseau Twitter/X} et fans endeuillé.es ont commencé à partager illégalement les films et séries et documentaires du réalisateur via des services comme Mega ou Google Drive. 
    
    \end{itemize}

Dans le temps qui nous a été imparti, nous avons également réalisé cinq entretiens. Notre but était de pouvoir fournir une diversité d'âge, de genre, de pratiques et d'expériences. Ces entretiens ont été pour la plupart réalisés en tête-à-tête, sauf pour Théo qui habite trop loin. Une transcription automatique a été systématiquement mise en place par captation audio et traitement utilisant l'outil Whisper d'OpenAI.\newline

Tableau des profils des différents entretiens.\\
\begin{tabular}{|p{2cm}|p{2cm}|p{2cm}|p{2cm}|p{2cm}|}

\hline
   Prénom  & Âge & Profession & Méthode & Résumé\\
\hline
    Adrien & 23 ans & Étudiant en informatique & Torrenting, Téléchargement & Diplôme en informatique \\
\hline
    Théo & 27 ans & Ingénieur informaticien & Torrenting, Serveur & Diplôme en informatique\\
\hline
    Hilaire & 26 ans & Étudiant en sociologie & Téléchargement, forums, Torrenting & Diplôme en communication et en sociologie \\
\hline 
    Zakary & 25 ans & Étudiant en communication & Téléchargement, communauté & Diplôme en histoire et communication \\
\hline
    Elise & 23 ans & Étudiante en communication & Téléchargement, communauté & licence langue et édition, master de communication \\
\hline
\end{tabular}\\

Nous avons tenté de représenter équitablement les genres, mais un temps de recherche insuffisant entraîne un biais persistant lié à l'âge. La plupart des répondants sont des personnes de ma génération (nées entre 1995 et 2003). Nous avons cependant réussi à mêler des profils de personnes proches de l'informatique (la majorité du corps social qui m'entoure) et d'autres qui piratent sans pouvoir se reposer sur des connaissances techniques aussi poussées. 
Malgré une cohorte limitée, la saturation sémantique des réponses de nos enquêtés a été relativement rapide. La plupart de nos répondants ont une pratique du piratage privilégiant l'accès le plus simple et le plus rapide aux ressources, homogénéisant rapidement les réponses aux questions. Les entretiens individuels sont difficilement des portes d'entrée vers les communautés patrimoniales pirates - qui sont par définition des œuvres collectives ; ils restent cependant très intéressants pour comprendre la manière dont des individus interagissent avec ces communautés sans pour autant en faire partie.

Nous avons également utilisé comme sources plusieurs bibliothèques pirates : Sci-Hub, LibGen et Zlibrary ; pour comprendre et étudier à la fois la culture pirate, les outils et l'expérience que les individus trouvent sur ces bibliothèques, ainsi que l'architecture et l'organisation de ces sites. 

\section{Méthodologie et problématique}

%méthodologie

Notre méthodologie s'inspire à la fois des \textit{fanstudies}, de la sociologie de la déviance et des études en archivistique et en sciences de l'information et des bibliothèques. Notre but est de rapprocher les expériences personnelles des personnes qui piratent et l'étude de l'architecture organisationnelle des bibliothèques pirates pour comprendre les éléments et les formes d'organisation dans les communautés patrimoniales pirates. Nous essaierons de comprendre et de cartographier les émotions patrimoniales appliquées à des objets culturels numériques, et la manière dont ces émotions sont le fondement de l'émergence d'une pratique archivistique populaire. Notre sujet d'observation sera les dépôts pirates, des bibliothèques et archives communautaires, autogestionnaires et clandestines. Nous étudierons à la fois l'aspect sensible et intime de l'émotion patrimoniale, du vécu pirate, mais également le contexte d'organisation et de remise en question des normes de distribution et de conservation de la culture sur Internet. Nous cherchons à explorer la culture et les références qui permettent à des individus qui ne se sont sans doute jamais vus de s'organiser, et d'impacter le circuit de la diffusion de l'information en cherchant à inventer des modèles alternatifs de gestion et de partage du patrimoine immatériel. \\

%annonce de plan 

D'abord, nous nous concentrerons sur la vie intérieure des pirates, de leurs expériences et l'apprentissage social qui les sépare des communautés patrimoniales pirates. Ensuite, nous aborderons la culture pirate, sa vision du monde et son lien avec l'émotion patrimoniale de sédition et les problématiques d'organisation qui peuvent en ressortir. Enfin, nous irons nous immerger dans la grande bibliothèque, en édifiant une typologie des différents types de dépôts pirates, de leurs organisations et des différents buts pour lesquels ils ont été créés.

%Nous nous essaierons dans ce travail de comprendre les dynamiques culturelles et les organisations sociales sous-jacentes aux pratiques des communautés patrimoniales pirates. 
%Comment les dynamiques culturelles et les organisations sociales permettent-elles l'émergence d'une pratique archivistique populaire, autogestionnaire et clandestine ? Notre recherche s'inscrit dans un contexte de remise en question international des normes de distribution et de conservation du savoir, posant de nouveaux défis aux institutions patrimoniales et aux acteurs entrepreneuriaux de la distribution culturelle. Nous cherchons à explorer la manière dont ces entités informelles s'organisent, pour comprendre l'impact sur les circuits de la diffusion de la culture, mais également pour envisager des modèles alternatifs de gestion et de partage du patrimoine immatériel.\\

%Nous étudierons d'abord le chemin individuel d'apprentissage, de partage et de mise en relation des individus avec les communautés patrimoniales pirates. Ensuite, nous aborderons la vision du monde des individus qui piratent et les liens entre émotions pirates et émotions patrimoniales. Enfin, nous irons visiter plusieurs dépôts numériques pirates marqués par les personnalités et intérêts des individus qui les composent. 



\chapter{De l'individu à la communauté patrimoniale pirate}

\begin{quote}
    \textit{L'autre inconvénient de ce type de bibliothèque est qu'elle permet, appelle, encourage la xérocivilisation. Cette civilisation de la photocopie, en dépit des grands avantages qu'elle procure, entraîne avec elle une série de dangers pour l'édition, même d'un point de vue légal. C'est avant tout la mort de la notion de droit d'auteur. [...] Sans compter qu'on peut emprunter le livre, le sortir et que certaines coopératives d'étudiants vous feront les photocopies sur des feuilles perforées prêtes à ranger dans des classeurs. Dans ces coopératives aussi on vous dit parfois qu'on ne peut pas photocopier un livre entier : j'ai eu ce problème avec certains de mes  étudiants. Ils me disent : \og Il nous faut trente copies de ce livre mais ils refusent \fg.[...] Je leur dis : \og Très bien, faites faire une photocopie, puis rapportez le livre à la bibliothèque ; après demandez vingt-neuf copies d'une photocopie. Il n'y a pas de droits sur une photocopie\fg.}\\

    Umberto Eco, \textit{De Bibliotheca}, traduit de l'italien par Eliane Deschamps-Pria, édition de l'Echoppe, 1986.
\end{quote}

Dans ce chapitre, nous nous concentrerons sur les origines du lien entre les individus, les pratiques du piratage, et les communautés patrimoniales pirates. Nous nous intéresserons d'abord à la naissance des carrières pirates, puis nous comparerons les différentes manières d'apprendre et de s'autonomiser chez les nouveaux pirates. Enfin, nous terminerons avec l'utilisation des compétences acquises par le piratage dans le monde professionnel et académique. Nous nous inscrirons ici dans l'héritage des travaux de Patrice Flichy, sociologue des sciences de l'information et de la communication ayant étudié la massification et l'utilisation d'internet chez les amateurs \footnote{Patrice Flichy, \textit{Le sacre de l'amateur, sociologie des passions ordinaire à l'ère numérique}, collection République des idées, éditions du Seuil, 2014}. La grande majorité de nos enquêtés sont des membres de la \textit{génération Z}\footnote{la génération Z comprend les personnes nées entre la fin des années 1990 et et le début des années 2010.} qui ont grandi avec le développement de l'informatique et de la télécommunication domestique grand public. L'étude de l'acquisition des compétences des pirates nous plongera également dans l'apparition massive de la sociabilisation en ligne.

\section{Ohé Moussaillons ! Dynamiques de découverte et de transmission chez les jeunes pirates}

Dans cette première partie, nous étudierons les moments de découverte du piratage chez nos enquêté.es, nous nous intéresserons spécifiquement à la dimension sociale et émotionnelle des premiers pas dans l'univers du piratage.

\subsection*{Définir la culture pirate}

Pour qu'un individu commence à pirater, il doit alors passer par une étape de découverte et d'acceptation d'un nouvel usage, et dans le cas du piratage d'une nouvelle culture. Même très jeunes, nos enquêté·es connaissent les limites légales de la copie, et savent que l'utilisation qu'ils en font est illégale. Ils se retrouvent donc face à un dilemme moral : respecter la loi ou choisir de l'enfreindre consciemment, et nous allons essayer de comprendre les éléments qui fondent leurs choix. Nous pouvons résumer la culture pirate en quelques clés \footnote{\cite{de_kosnik_rogue_2016} page 2}.

\begin{enumerate}
    \item La culture pirate est l'assomption que tout est possible sur le réseau, de la citation, la transformation et le partage de biens culturels, sans citer et sans en demander l'autorisation à l'auteur. Ce cadre de pensée est également la base d'objets culturels comme les \textit{les memes} dont le principe est de tordre et réinventer un bout d'objet culturel sorti de son contexte à l'infini, les \textit{fanfictions} qui prennent un univers existant pour le transformer selon les envies des auteurices, ou les \textit{mods} dont le but est d'ajouter des fonctionnalités à un jeu, quitte à le transformer radicalement.
    \item Le piratage n'est pas assimilable à du vol, ni à un détournement de propriété, mais bien à la production non autorisée d'une copie numérique d'un contenu culturel. De ce fait, il existe une différence entre le cadre légal et l'éthique, et l'un ne définit pas l'autre\footnote{Richard Stalman,
\textit{Annexe 4 : Quelques expressions trompeuses qu'il vaut mieux éviter d'employer}, dans O. Blondeau, Libres enfants du savoir numérique : Une anthologie du “Libre” (p. 359-364). Éditions de l'Éclat, définition de vol.}
\end{enumerate}

La découverte des pratiques de piratage de biens culturels est un parcours relativement homogène chez mes enquêtés. Sauf quelques rares exceptions, les individus découvrent le piratage avec leurs premiers équipements numériques (premier téléphone, premier mp3, ordinateur familial, premier pc). Cette découverte passe par un rapport entre adolescents qui ont les mêmes équipements, ou par une personne extérieure (parents, amis des parents, grands frères et sœurs) que l'on peut surnommer comme étant des \textit{figures tutélaire}. L'informatique devient alors l'interface et le principal médium de l'intérêt pour la musique, le cinéma, les séries, les jeux vidéo et les mangas.

\subsection*{La découverte par les pairs}

La voie majoritaire de découverte du piratage se passe dans le cadre amical durant l'enseignement secondaire. En discutant avec ses ami.es de contenus culturels, principalement de musique, l'individu prend conscience de l'existence d'outils et de communautés dédiés à l'acquisition pirate de ses centres d'intérêts. Ce chemin d'apprentissage privilégie l’acquisition de compétences de surface, permettant l'acquisition de contenus en passant par des services tiers. Youtubemp3, un site de conversion de vidéo YouTube en fichiers mp3 aujourd'hui disparu, mais dont des milliers de clones ont pris la relève, a été pour la majorité de mes enquêtés la première étape de la découverte du piratage. La principale méthode de diffusion de ce site est alors la discussion entre ami.es durant la récréation du collège. La découverte de l'outil creuse ensuite la curiosité au fur et à mesure que les goûts culturels de l'individu s'autonomisent et que l'informatique se démystifie petit à petit.

Hilaire et Élise sont deux profils ayant commencé à pirater de cette manière. Hilaire s'est formé au piratage avec ses groupes de fans et de partage de dramas coréens traduits. Le cœur de l'activité de pirate d’Élise réside dans le téléchargement de livres, essentiellement des classiques, des livres qu'elle possède déjà, ou des livres de romance qu'elle n'oserait pas acheter en magasin.

\begin{quote}
H - \textit{Ma grand-mère m'avait offert un ordi portable qui coûtait très cher à l'époque.
 [...] Au début, je pense que ça a commencé par la musique, je pense que c'est le tout premier truc. C'est genre les sites où tu mets le lien ... Youtubemp3 un truc comme ça. [...] J'ai encore tous mes fichiers, donc ça, c'est un peu marrant, je pourrais aller voir les tout premiers trucs que j'ai enregistrés.}
\end{quote}

\begin{quote}
    E - \textit{Déjà, quand j'ai commencé à pirater, j'avais un ami. Au départ ce n'était pas vraiment moi qui le faisais. J'ai le père d'une amie proche qui, lui, a téléchargé énormément de films et tout, et il me les passait sur clé USB. }
\end{quote}

La découverte du piratage entre pairs favorise une découverte graduelle et relativement lente de solutions peu techniques. Ici, la conceptualisation de l'apprentissage du piratage ressemble à une découverte géographique plutôt qu'à des compétences informatiques. Les individus se partagent des \og bonnes adresses \fg et nous pouvons résumer leur horizon mental à une carte s'enrichissant avec le temps de lieux (des sites internet). Cependant, cette capacité de recherche de solutions fait que ces gens ont en général un niveau satisfaisant de maîtrise de la technologie \footnote{Unesco Regional Bureau for Education in Asian and Pacific \textit{Diverse approaches to developing and implementing competency-based ICT training for teachers: a case study}, TH/EISD/16/019-500, 2016}, même si cela n'est que le médium entre eux et leurs passions. Ces gens entrent rarement en contact avec des communautés patrimoniales pirates, sauf pour le cas particulier du piratage académique dont nous allons reparler un peu plus loin.\\

%A compléter et enrichir
Dans les archives des skyblogs de la BnF, nous trouvons plusieurs blogs dont le but est de partager des mangas ou des films que l'on a appréciés, tout en partageant également un lieu de téléchargement ou de streaming pour ses ami.es et followers. Par exemple, le blog de Spymon : \textit{telechargement09} \footnote{voir Annexes Skyblogs} qui se concentre sur le partage de sa passion pour la musique et l'esthétique du genre métal/gothique. Nous pouvons remarquer que les personnes téléchargeant ou partageant de la musique au format mp3 sur internet ont souvent le réflexe de ne pas assumer cette action comme du piratage. C'est également produit dans un de nos entretiens où Élise, cherchant dans sa mémoire l'instant où elle a commencé à pirater, s'est rendu compte au fur et à mesure des exemples que sa première rencontre avec la piraterie était le téléchargement de musique. 

\subsection*{La découverte par une figure tutélaire}

La deuxième manière dont on peut découvrir le piratage est en passant par une \textit{figure tutélaire}, une personne déjà très à l'aise avec l'informatique, la technologie et le piratage, qui va prendre le rôle de messager vertical entre les communautés patrimoniales pirates et l'individu. Cette manière de découvrir le piratage peut imprimer deux manières de faire chez les personnes qui commencent à construire une habitude de piratage. La première, sans doute la plus commune, est un apprentissage plus complet et une montée en compétences très rapide de ces individus. On y retrouve des pratiques comme le partage en pair à pair, le torrent, le partage d'applications modifiées, la mise en place de serveurs, de seedbox, ou de VPN comme Hamachi pour faciliter le partage ou le lien avec ses ami.es. L'exemple type est celui d'Adrien, qui a été introduit aux pratiques de piratage très jeune grâce à son père qui lui apprend et lui finance des moyens de pirater dans une sécurité relative.

\begin{quote}
    A - \textit{Donc depuis tout petit, moi je vois des ordinateurs à la maison, [...] avec ma mère j'ai découvert l'informatique et la bureautique, et avec mon père plus tard j'ai découvert le reste. Donc beaucoup de jeux vidéos, et aussi de la consommation de contenu, illégal particulièrement. C'est mon père qui m'a fait découvrir le monde merveilleux de la R4\footnote{La R4, pour \textit{Revolution for DS} est un \textit{linker}, une cartouche permettant d'ajouter énormément de fonctionnalités à une console, en particulier tricher, ou faire fonctionner des ROMs de jeux sans avoir de cartouche officielle. S'il n'existe pas de loi encadrant l'existence des linkers et leur utilisation, des revendeurs ont été condamnés en France, voir l'article de Julien Lausson, \textit{La vente de linkers pour Nintendo DS jugée illégale en France}, Numérama, 04/10/2011} La vente de linkers pour Nintendo DS jugée illégale en France  à l'époque, et surtout le site EmulIsland, je crois que ça s'appelait, pour les jeux Wii, je m'en souviens encore. Et donc ça c'était avec mon père que j'ai découvert ça.}
\end{quote}

Une autre évolution possible de la découverte par une figure tutélaire peut être le fait de continuer de se reposer sur l'expérience de cette personne. Ceci semble arriver chez des gens qui n'ont pas l'envie, le temps, ou la curiosité d'approfondir leur rapport à l'informatique. C'est par exemple le cas de Zakary avec son frère. Les deux ont toujours été très proches malgré la différence d'âge et l'un des principaux vecteurs de cette complicité fraternelle est l'amour de la  musique et du cinéma d'animation japonais.

\begin{quote}
    Z - \textit{J'ai un grand frère [...] et c'est cette personne-là qui m'a introduit au piratage. Quand j'étais très jeune, un peu avant le collège, il m'a un peu expliqué sur l'ordinateur familial comment récupérer les vidéos YouTube que j'aimais bien, comment mettre des musiques sur mon téléphone portable pour que je puisse les écouter quand je voulais. [...] Mon frère avait certains intérêts que moi, je n'avais pas forcément. Et donc, quand je voulais ouvrir une porte vers ces mondes-là, je n'avais pas forcément les codes. Je savais que j'avais un peu son héritage qui pouvait être là. J'ai même bien parlé avec lui pour qu'il me donne quelques clés.}
\end{quote}

L'un des points culminants de cette relation trouve son incarnation dans un objet, un iPod rempli de ses musiques piratées, que le frère de Zakary lui offre. Ce qui semble expliquer que Zakary n'ait pas eu le même parcours qu'Adrien est qu'il est plus intéressé par les goûts musicaux et cinématographiques de son frère que par les méthodes et la \textit{sociabilisation pirate}.
    \begin{quote}
        Z - \textit{Cela a définit mes goûts musicaux pour le restant de mes jours.}
    \end{quote}

\subsection*{La découverte en autonomie}

Une troisième manière de découvrir le piratage nous est indiquée dans le vécu de Théo, qui a grandi dans un lieu isolé en pleine campagne, avec l'accès à un ordinateur dès la classe de sixième. Avec des parents qui utilisaient un minitel pour des tâches administratives, la famille de Théo est une famille qui connaît et utilise des services informatiques, sans pour autant avoir une figure tutélaire technophile. La spécificité du parcours de Théo est l'autoformation, conséquence immédiate de sa passion naissante pour l'informatique qu'il partage avec son frère, qui lui offre la connexion à un monde culturel qui lui manque.

\begin{quote}
    T - \textit{J'étais assez isolé en soi des grandes villes et tout [...] j'ai assez rapidement eu accès à un ordinateur, comment utiliser un ordinateur, comment utiliser internet, des cet âge là. Par contre mon téléphone portable je ne l'ai eu qu'à partir de la seconde. Je n'avais pas le droit avant, ça peut paraître un peu bizarre d'ailleurs maintenant sur l'évolution d'internet. [...] J'ai un peu trouvé ma passion là dessus sur l'informatique de manière générale ...}\\

    \textit{Mon utilisation au tout début c'était au lycée où je commençais à vouloir trouver des films. Je reviens au fait que je vivais en campagne mais en fait j'avais pas mon autonomie, et je pouvais pas aller au lycée, je pouvais pas aller au cinéma à pied ou en bus ou quoi. Donc ça dépendait de mes parents, mes parents n'étaient pas très intéressés par le cinéma, moi j'adorais ça. Du coup le but c'était d'aller chercher ce que je pouvais sur internet et donc assez rapidement je suis tombé sur des méthodes de téléchargement, comment essayer de m'en sortir tout seul ...}
\end{quote}

C'est cette combinaison de l'informatique comme un médium pour une passion pour le cinéma et une passion pour l'informatique qui permet à Théo de briser la barrière entre son isolement et les communautés pirates. \\


\section{Grandir en piratant et évolution des pratiques au fil de la vie}

\begin{quote}
    \textit{La démocratisation des compétences repose d'abord sur l'accroissement du niveau moyen de connaissances (dû notamment à l'allongement de la scolarité) et sur la possibilité offerte par Internet de faire circuler les savoirs, de livrer son opinion à un public plus vaste. L'amateur qui apparaît aujourd'hui à la faveur des techniques numériques y ajoute la volonté d'acquérir et d'améliorer des compétences dans tel ou tel domaine.}\\

Patrice Flichy, \textit{Le sacre de l'amateur, sociologie des passions ordinaires à l'air numérique}, 2014
\end{quote}

L'un des éléments les plus intéressants dans l'étude des carrières pirates est l'évolution de la pratique au cours de la vie de nos enquêté.es. En devenant de jeunes adultes, les individus se découvrent de nouvelles passions, et, entrant dans les études supérieures (tous mes enquêté.es ont fait des études supérieures jusqu'à minimum bac+3), iels réutilisent ce qu'ils faisaient pour les films et les livres dans le cadre académique. Nous nous concentrerons dans cette partie sur l'évolution et l'adaptation des individus qui ont appris à pirater durant leur adolescence et qui vont continuer leur carrière pirate dans leurs études, leur vie professionnelle ou personnelle. Nous déshabillerons leurs routines, leurs justifications, et les évolutions qu'iels ont dû mettre en place pour s'adapter à un contexte qui change. % reformuler la dernière phrase de merde

\subsection*{Pérennisation, transformation ou abandon des habitudes pirates des jeunes adultes}

La vie de jeune adulte semble produire deux conséquences sur les habitudes de piratage de nos enquêté.es. Certains, comme Théo et Adrien, semblent creuser et intensifier leurs habitudes de piratage ; d'autres, comme Hilaire et Élise, semblent mettre en pause pendant un temps leurs carrières de pirates. Ces moments de pause semblent liés à l'émergence de nouveaux types de produits culturels comme la vidéo à la demande et les abonnements de streaming dont le représentant principal est Netflix. La praticité et le coût très faible ont permis à certaines entreprises, comme Netflix, Spotify ou Steam, de remplacer les habitudes des pirates par des alternatives légales. Le changement de rythme de vie, de nouveaux temps de transport, le fait d'avoir plus de moyens à mettre dans sa consommation culturelle, sont diverses raisons d'abandon des carrières de pirates des jeunes adultes. Le téléchargement demande du temps et une organisation qui ne peuvent pas être compatibles avec la vie d'une personne s'insérant professionnellement ou qui habite loin.  

\begin{quote}
    H - \textit{J'ai un abonnement Netflix, et je pense que ouais du coup à cette période j'arrête totalement de pirater films, séries, j'ai un peu continué pour de la musique durant ma licence, mais ça n'a pas duré très longtemps parce qu'après j'ai eu un compte Spotify. Et c'était quand même vachement pratique de ne pas avoir à aller chercher tous les trucs que tu as envie d'écouter, de les télécharger et de mettre sur ton téléphone. À la limite je pense que j'ai dû aller quelques fois sur des sites pour lire des scans de mangas. Mais ça n'a pas été vraiment très très dur, mais ça m'est pas arrivé beaucoup. }
\end{quote}

L'évolution d'une pratique de piratage peut signifier l'arrêt ou le ralentissement d'une partie de l'activité d'un individu, sans pour autant signifier l'essoufflement total de la carrière de pirate. Par exemple, Élise pirate énormément de livres, et continue de pirater des livres, mais elle a arrêté de télécharger des livres audio, auxquels elle accède maintenant par un service légal qu'elle paie. 

\begin{quote}
    E - \textit{Moi, c'est venu vraiment après, quand j'ai découvert le monde merveilleux du piratage de livres. Même des livres audio, des fois, j'en prenais aussi. Ça, d'ailleurs, c'est un autre moyen puisque les livres audio, avant, je les piratais. Mais maintenant il y en a plein sur Spotify, par exemple. Donc,il y a même des choses que je pensais illégales. Maintenant je les fais ... j'ai un abonnement Spotify chaque mois. Mais je peux lire, je peux écouter les Hauts de Hurle-Vents, je sais pas, enfin gratuit quoi, enfin \og gratuit \fg entre guillemets. Mais ça c'est venu récemment, je crois.}
\end{quote}

Le sujet animant toute la communauté travaillant sur le piratage est l'impact d'offres légales jugées supérieures à des expériences pirates \footnote{Anthony Galluzo à toute une sous partie sur cette question de l'impacte des offres légales de streaming sur les comportements illégaux.}. Cependant, cette question semble beaucoup plus sensible au niveau de revenu des individus que l'existence intrinsèque d'offres légales. Par exemple, avant l'apparition de Netflix, il existait durant les années 1990-2000 des magasins de prêt-vidéos permettant un accès légal à des films pour une fraction du prix de l'achat d'un DVD. L'apparition des plateformes de streaming n'a pas révolutionné le cœur de la question du piratage, mais elles ont cependant modifié les échelles auxquelles ces questions se posent. \\

Hilaire a un parcours particulier, après avoir presque totalement arrêté toute forme de piratage durant ses études supérieures et le début de sa vie professionnelle, il recommence à pirater durant sa reprise d'études d'un master recherche en sociologie. Sa reprise du piratage est ici également directement liée à l'impact de sa reprise d'études sur ses revenus. 

\begin{quote}
    H - \textit{Bha depuis ce temps bah du coup bon ça je te l'apprends pas mais au fil des années les plateformes Netflix, Spotify augmentent leur prix. Donc à un moment où moi ça m'a saoulé, du coup j'ai juste arrêté d'utiliser Netflix et Spotify? D'abord Netflix parce qu'en vrai c'était plus facile d'accéder à du streaming de séries et de films que je trouvais, et du coup la transition était assez facile, de revenir au piratage notamment par du torrent.[...] D'abord je suis passé sur Deezer parce que c'était ... enfin c'est une entreprise française et ils rémunéraient un peu mieux les artistes. Du coup j'ai d'abord fait ce choix en mode ... enfin un choix éthique quoi. Et après j'ai arrêté parce que je me suis rendu compte qu'en fait j'utilisais pas tant que ça Deezer. [...] Et récemment j'ai un ami qui m'a envoyé le lien pour télécharger une APK de Spotify et avoir Spotify sans pub, sans payer quoi. Et du coup j'ai un peu utilisé mais là ça commence à buger du coup il faut que je le retélécharge dans d'autres version pour voir si ça remarche. Donc je n'ai pas encore trouvé de solution pour la musique.}
\end{quote}

Nous reviendrons sur le point de jonction entre le vécu économique des individus et les arguments éthiques et politiques qu'ils déploient dans le chapitre suivant. Cependant, il est intéressant d'évaluer l'argument que nous avions déjà rencontré dans notre introduction, celui des pirates comme une manne de consommateurs potentiels. S'il n'est pas faux que des services comme Netflix, Spotify ou Steam aient réussi à générer un marché qui a converti une partie des populations pirates en clients, et que de nombreux spécialistes prédisaient la mort de la pratique du piratage ; la multiplication des plateformes et donc des abonnements, l'exclusivité du contenu, l'accès direct à la vidéo à la demande a redonné un souffle à la pratique. Les manières de pirater ont évolué en tenant compte de ces nouveaux acteurs et de l'apparition de l'économie de la publicité. En tapant le mot \og Torrent \fg dans l'outil Google Trends, nous possédons un témoin relativement fiable de l'intérêt du public pour le piratage utilisant des liens magnétiques depuis 2004 à nos jours \footnote{voir annexes Résultats Google Trends pour le mot Torrent}. Nous comprenons que la popularité du torrent s'effondre à partir de la deuxième moitié des années 2010 ; cependant, cela n'indique pas que les individus piratent moins mais que les méthodes de piratage évoluent. Une nouvelle manière de consommer du contenu pirate est le streaming, à la manière de Netflix ou Spotify. Les hébergeurs de contenus piratés ont adapté dans la grande majorité leur fonctionnement pour suivre les habitudes des consommateurs, abaisser les barrières à l'entrée (plus besoin de maîtriser le torrenting, plus besoin de télécharger un fichier), et rassurer les utilisateurs effrayés par l'association entre protocole Torrent et la surveillance d'Hadopi. Si Netflix et Spotify étaient des services révolutionnaires et peu onéreux pendant une époque, l'infrastructure économique et la superstructure culturelle de la société ont transformé cette pratique en une norme, et le fait social pirate, comme toutes les pratiques de déviance, a muté pour s'y insérer. 

\subsection*{La piraterie académique}

Dans les différentes manières dont les individus réutilisent, adaptent ou abandonnent leurs pratiques pirates, le piratage académique trouve une place particulière qui symbolise souvent le passage de l'adolescence (pirater pour de la consommation culturelle ludique) à la vie adulte et professionnelle (pirater pour accéder à de la connaissance et des compétences).

\begin{quote}
    \textit{A - Et donc ouais la prépa. Ça c'est un peu calmé. Je fais mes deux ans de classe préparatoire. Il y a quand même du contenu que j'ai téléchargé, qui est surtout en fait marginal, on avait la chance d'avoir une bibliothèque qui était bien fournie, donc on avait des manuels, mais il y avait certains manuels qu'on ne pouvait pas avoir autrement qu'en téléchargeant. Donc là, ça a été mon premier contact avec le téléchargement académique . Même si pas encore, parce que je ne savais pas tout à fait ce que c'était encore les articles de recherche vraiment, mais tout ce qui va être manuels d'histoire, manuel de maths, des trucs comme ça. Je ne pouvais pas lâcher 50 balles dans un manuel, et tu avais trois manuels par matière, ce n'était pas envisageable. Donc là, j'ai commencé à trouver des trucs sur internet. C'était bien.}
\end{quote}

Si le piratage académique accompagne le passage au monde adulte et professionnel, il accompagne donc également l'éveil des sensibilités politiques et des émotions patrimoniales. Les enquêtés ayant mentionné un rapport politique, une vision du monde ou une considération patrimoniale liée à la pratique du piratage l'ont intégré à leurs récits au moment de l'acquisition d'autonomie et de maturité des études supérieures et de l'entrée dans le monde professionnel. 

\begin{quote}
    T- \textit{C'est pas anodin effectivement [...] J'ai l'approche très scientifique de cette chose-là. C'est que pour moi, la science, en tout cas, doit être entièrement ouverte. J'ai beaucoup de mal avec le fait que les labos s'enferment sur des articles qui ne sont publiés qu'entre eux, que certains industriels ferment certains verrous sur l'accès à la science et à des technologies de pointes. Moi je suis pour, effectivement le partage complet, peu et simple, en fait, en science. Et du coup ça transpire un peu aussi sur ma conception du partage culturel. Je suis un gros soutien pour le fait que l'accès à la culture, en particulier pour les jeunes, doit être le plus simple possible. Donc ... en particulier sans frais, parce que pour les jeunes, forcément, le premier frein, ça doit être le budget. Donc le fait d'encourager les jeunes, que ce soit aller au musée, que ce soit aller au cinéma, que ce soit aller à une bibliothèque ou quoi, pour moi, tout doit être le plus simple possible.}
\end{quote}

Cette citation de Théo, après une demi-heure d'entretien, résonne avec des discours qui nous sont familiers, nous pouvons citer ici la page \textit{About} du site de médiation artistique pirate Ubuweb \footnote{Voir l'annexe Message d'introduction à Ubuweb}, ou avec le \textit{Manifeste pour une guérilla pour le libre accès} d'Aaron Swartz, ou du slogan du site Sci-Hub \footnote{« To remove all barriers in the way of science » (« Éliminer tout obstacle sur la voie de la science »)}. La plupart de nos enquêté·es ne se définissent pas comme des personnes particulièrement politisées. Si iels ont des idées et des sensibilités politiques, iels ne sont ni encarté·es, ni militant·es dans des organisations politiques. Cependant, la pratique du piratage amène un discours qui utilise un vocabulaire et des idées libertaires. Si ces personnes n'ont pas le temps et les moyens d'acquérir les compétences pour aligner leurs logiques et leurs actions, sociabiliser comme \og pirate \fg ou amateur, ou n'importe quel autre terme demande de rentrer en contact avec des communautés patrimoniales pirates, et donc de comprendre et de faire sienne une partie de la culture pirate et de ses influences comme la culture du libre et les idéologies libertaires. 

\begin{quote}
    H - \textit{Je ne me revendique pas de ça mais parce que le sujet ne me vient pas quoi. Enfin c'est pas trop quelque chose dont les gens autour de moi parlent. [...] Mais ouai je suis pas sûre que ce soit vraiment une identité même si en fait ... Ça fait depuis ... C'est en licence que j'ai découvert la culture du libre. Et j'avais eu des cours sur ça en fait par un intervenant qui travaillait dans un Fab Lab. Et je pense que vraiment ce cours a vraiment changé toute ma perspective de voir le monde. J'étais en mode ah mais oui c'est génial en fait. Genre tu peux hacker des machines, mais tu peux aussi hacker des vêtements, des trucs, enfin tu peux tout faire toi-même. [...] Ouai il nous a grave initié à la culture du libre. Wikipédia, Open Office, tout ça [...] Du coup c'est à cette époque que j'ai découvert un peu que c'était un mouvement de pirater des trucs, [...] j'avais des pratiques de piratage pas je savais pas que ça pouvait s'inscrire dans un mouvement, et dans des revendications politiques.}
\end{quote} 


\subsection*{De pirates à corsaires :  Quand le piratage s'invite dans le vie de l'entreprise}

Précédemment, nous avons exploré l'évolution des compétences de piratage dans le domaine académique et des études supérieures, mais cette mécanique s’étend au-delà des études. Pour certains, le piratage devient une facette de leur insertion professionnelle, où des entreprises recherchent et les rémunèrent pour leurs compétences. Certains milieux professionnels et secteurs économiques jouent avec les limites dans l'utilisation de certaines utilisations de contenu licencié. En discutant avec le comité de sélection \textit{L'Association des Archivistes Français}, une membre du jury me confie que le sujet de ce mémoire résonne particulièrement avec son vécu, les publications des normes AFNOR sont trop onéreuses pour son service d'archives, et son équipe a décidé de pirater ces documents pour pouvoir exercer correctement son travail. Cette même archiviste me parlera ensuite de l'affaire Filaé \footnote{Pierre Alexandre Compte, \textit{Archives : « L’affaire Filae », du besoin de repenser la politique de diffusion des données culturelles} dans la Gazette des communes, 21/12/2016}, anciennement Généalogie.com, une entreprise de généalogie ayant réalisé du profit sur la mise à disposition de documents numérisés et partagés gratuitement par les départements. J'aimerais introduire ici le mot \textit{bureautisation}, qui décrit à la fois l'économie de service tertiarisé et l'agencement physique des moyens de production au sein de l'espace du bureau, mais aussi de l'interface bureautique et donc des logiciels qui se font l'interface entre un.e travailleureuse et sa mission. La bureautisation amène des individus à privilégier l'alternative pirate : pour ne pas avoir à couper dans leurs budgets, pour ne pas avoir à faire une demande administrative chronophage à la hiérarchie, pour se tenir à jour des pratiques professionnelles de leurs branches.
Dans les milieux professionnels demandant un va-et-vient constant avec la recherche (milieu académique, professorat, ingénierie, médecine), le piratage devient le seul moyen de pouvoir accéder aux ressources pour que les étudiants, les apprentis et les professionnels puissent mettre à jour et progresser dans leurs compétences.

\begin{quote}
    Sylvie\footnote{Sylvie ne fait pas partie de nos enquêtés pour la simple et bonne raison qu'elle ne pirate pas, je la cite simplement suite à une discussion intéressante sur la place du piratage dans le monde des professions médicales} - \textit{Je donne des cours de français médical à des jeunes médecins immigrés, et je suis impressioné. Au début de l'année ils m'ont tous demandé le titre du manuel que j'utilise, et ils l'ont tous piraté, téléchargé, et ils se sont refilés les bonnes combines, c'est très impressionnant.}
\end{quote}

Le monde de l'intelligence artificielle, spécifiquement les grands modèles de langage (LLM) captent l'attention, les investissements et la couverture médiatique. Ces modèles demandant énormément de données en entrée, c'est logiquement que plusieurs entreprises ont été prises en train d'avoir recours au piratage de contenus culturels et scientifiques. L'exemple le plus récent est celui de la corporation de Mark Zuckerberg, Meta, un des acteurs les plus importants dans l’entraînement et la mise à disposition de LLM au grand public, qui a vu ses mails internes divulgués en démontrant un usage systématique de contenu licencié \footnote{Ashley Belanger, \textit{“Torrenting from a corporate laptop doesn’t feel right”: Meta emails unsealed}, sur le site arstechnica.com, 06/02/2025}. Meta s'est contenté de télécharger plus de 70 téraoctets de livres et papiers numérisés sur la plus grande bibliothèque pirate existante, Libgen. Les exemples d'acteurs privés épinglés pour du piratage culturel, ou du scraping de données personnelles dans le but de former des bases de données d’entraînement d'intelligence artificielle demanderont un mémoire à part entière pour être traités. 

\begin{quote}
    T - \textit{C'est compliqué à dire mais je l'ai même fait à la demande de mon employer assez étrangement, mais en fait c'est quelque chose qui est aussi répandu dans le milieu scientifique en particulier pour ce qui est de l'intelligence artificielle, le cœur du sujet, la base du travail dans l'intelligence artificielle c'est les bases de données. Si tu n'as pas de données tu ne peux pas entraîner de modèle donc tu ne peux rien faire et le problème c'est que les bases de données il faut les constituer. Il faut avoir des données propres, des données qui font du sens et du coup pour ça il faut avoir des données qui viennent de quelque part. Si ton entreprise n'est pas capable de générer elle-même les données qui l’intéresse il faut bien qu'elle aille les chercher ailleurs pour entraîner ses modèles et donc du coup très souvent tu vas avoir recours à du vol de données pur et simple. [...] Certains scientifiques et certaines entreprises ne se gênent pas pour aller voler ces jeux de données s'ils le peuvent [...] c'est que pour mon stage de fin d'études j'étais employé au sein de l'entreprise qui m'ont demandé expressément d'utiliser mon compte étudiant pour faire des demandes à des chercheurs qui sinon auraient fait payer les entreprises pour accéder à des données. J'ai eu aussi le cas de simplement aller essayer de piller des jeux de données sur des sites par forcément recommandables mais qui avait réussi à partager les jeux de données qui étaient utilisés pour différents articles.}
\end{quote}

\section{Conclusion}

La découverte du piratage se fait principalement par trois voies : la socialisation par les pairs, l'influence d'une figure tutélaire, et la curiosité personnelle. Chaque expérience apporte une perspective unique et offre des capacités d'apprentissage différentes aux individus. Il est important de préciser que ces différentes catégories s’interpénètrent et se superposent selon les situations. Ces interactions leur permettent d'acquérir des outils essentiels pour leur parcours académique ou professionnel, soulignant l'importance de la socialisation et de l'apprentissage continu dans le piratage. Comme le souligne Anthony Galluzo, la pratique du piratage est un ratio entre la capacité d'apprentissage, la mise en place de routines d'accès au contenu piraté, et l'entretien de ces compétences et routines. Les individus peuvent alors continuer ou reprendre leurs habitudes pirates et adapter les compétences qu’ils ont développées pour avoir accès à des ressources culturelles ou logicielles. Ce transfert de compétence se met en place durant les études supérieures et le besoin d’avoir accès à de la littérature scientifique, et peut continuer dans le monde professionnel en fonction de si la personne a besoin de continuer d’avoir accès à de la littérature scientifique, de la documentation technique, ou des outils logiciels et bases de données.

 
%III/ Dépôts numériques, hybride entre un lieu de
%conservation et de partage.
%• Le dépôt pirate, entre archive et bibliothèques
%– La fonction de conservation des dépôts numériques.
%– Le partage et la diffusion des ressources piratées.
%• Les communautés autour des dépôts numériques
%– Les contributeurs et les utilisateurs des dépôts.
%¨ Les dynamiques sociales et les interactions au sein des communautés.
%2
%– Les contributeurs et les utilisateurs des dépôts.

\chapter{Ni dieux, ni maîtres, ni DRM}

\begin{quote}
    \textit{ Ceci n’est ni un parti, ni un programme, ni un slogan, ni une méthode, ni quoi que ce soit d’autre qu’une idée à développer. Postulat : Chaque action sur le monde, de la plus petite à la plus grande, mérite d’être observée méticuleusement et ses conséquences envisagées à l’échelle globale, lointaine, holistique. }\\

    Anonyme, \textit{Manifeste pour une révolution liquide}, 2019
\end{quote}


Dans ce deuxième chapitre, nous étudierons la perception du monde des individus qui composent les communautés patrimoniales pirates. Nous nous intéresserons d'abord aux sentiments de sédition et de rébellion des individus et des communautés face à l'ordre des choses. Ensuite, nous nous concentrerons sur les émotions de peur, d'inquiétude et d'incertitude face aux abysses numériques. Enfin, nous creuserons les parallèles entre la vie culturelle en ligne et les engagements personnels dans le monde physique. Nous nous inscrirons ici dans le travail sur les émotions patrimoniales de Daniel Fabre \footnote{\cite{fabre_emotions_2013}} et les travaux qui enrichissent ce travail, comme ceux d'Emmanuelle Bermès en l'introduisant dans des contextes de communautés patrimoniales numériques \footnote{\cite{bermes_lecran_2024}}. Nous nous écarterons également du pur travail de restitution d'entretiens que nous avons fournis dans le chapitre précédent. Nous nous plongerons dans les traces de plusieurs communautés patrimoniales afin d'affiner ce lien entre émotions individuelles et ces communautés.

\section{Les émotions de la clandestinité}

L'une des plateformes les plus intéressantes à explorer et observer est sans conteste Reddit. Surnommé  \textit{Le cœur de l'Internet}, ce site d'agrégation de forums permet de trouver des communautés dédiées à n'importe quel sujet. Concernant la piraterie culturelle, Reddit et ses forums dédiés sont d'une accessibilité, d'une transparence et d'une documentation rares et précieuses pour le milieu et ses codes. 

\subsection*{La peur, l'inquiétude et la prise de risque}

Dans ce monde, la transparence et une documentation claire sont l'exception qui confirme la règle, la plupart des services d'accès au piratage culturel demandent un minimum de compétences et des savoirs prérequis. Cette opacité totale, doublée de l'inquiétude de télécharger des fichiers caviardés, d'être sur les radars des autorités sanctionnant le piratage, fait que le vécu des pirates ressemble bien plus à une traversée dans l'inconnu qu'à une chasse au trésor. Dans les verbatims de nos enquêtés, l'émotion présente à chaque étape du processus d'exploration, d'acquisition de compétences et de découverte de communautés pirates est la peur ou la crainte. 

\begin{quote}
    A - \textit{Je pense qu'il y a eu une partie de peur au début, quand c'est moi qui m'y suis mis à le faire. Quand mon père le faisait, je regardais derrière son écran, et ça avait l'air trop bien et super facile. Quand c'est moi qui ai dû aller sur des sites un peu shady, avec les extensions, c'est pas que .fr, .com, .org, mais des trucs un peu bizarres, qui sont en fait normal mais sur le coup t'es pas habitué, comment ça .xyz ? C'est quoi cette merde .me ? [...] Je pense que la première étape c'était un peu la peur, un peu l'excitation de ses découvertes [...] de se rendre compte que t'as une sorte de monde, enfin pas un monde caché ce serait un peu abusé, mais de découvrir au fur et à mesure tout ce qui possible de faire avec pas grand chose.}
\end{quote}

\begin{quote}
    T - \textit{ J'avais un fond de crainte et je n'étais pas complètement renseigné non plus, mais j'ai toujours été ce qu'on appelle le mauvais utilisateur dans le torrent c'est celui qui vient se servir et qui n'aide personne [...] parce que c'est comme ça que tu te fais avoir, c'est quand toi tu émets des fichiers, la plupart du temps pour le torrent. [...] Je triche, parce que je fonctionne aussi avec des sites qui demandent à ce que tu équilibres ton ratio, mais tu as des manières de tricher, t'as des manières de fausser l'envoi de certains fichiers et ça compense à chaque fois ce que tu télécharges. [...] Je jour contre les  règles, mais la logique derrière c'est pas d'être le vilain de l'histoire, c'est de ne pas me faire avoir.}
\end{quote}

La source principale d'inquiétude et de peur chez les pirates sont les autorités judiciaires. Chez nos enquêtés, la période d'activité d' Hadopi a laissé un mélange de crainte et de défiance. La majorité d'entre eux a reçu ou a connu une personne ayant reçu une des lettres et des mails de la haute autorité. Si aujourd'hui les lettres Hadopi sont une histoire rigolote pour les pirates francophones, les mouvements étatiques pour le contrôle et la surveillance du web, en particulier pour le piratage, ont produit et contribué à alimenter de réelles inquiétudes dans les communautés. 

\begin{quote}

A - \textit{C'est aussi le moment où j'ai découvert le mot-clé Hadopi, quand j'ai reçu une lettre d'Hadopi, en fait, chez moi. [...] Ouais, j'ai reçu une lettre Hadopi, mais ouais je m'en souviens, c'était un 31 octobre je crois la date même. Alors la question c'était est-ce que c'était moi ou est-ce que c'était le beau-père de l'époque ? On ne sait pas, mais l'adresse IP a été fag, on a reçu une lettre d'Hadopi, et c'est à ce moment-là que mon père m'a payé l'abonnement à la seedbox, comme ça je pouvais télécharger autant que je voulais à l'extérieur, et le principe des torrents étant quand même de télécharger pour rendre ensuite. La seedbox était constamment ouverte, donc sur un serveur, que je pouvais seed constamment, les trucs que je téléchargeais pour pouvoir les envoyer aux autres.}
    
\end{quote}

\subsection*{Trouver et utiliser un service clandestin}

Au-delà de la réflexion tautologique  \og les communautés pirates sont difficiles à découvrir car elles sont illégales et donc se cachent \fg, il est intéressant d'essayer de comprendre le fonctionnement et les organisations mises en place par ces communautés pour survivre et mettre en place des services communautaires. Dans notre introduction, nous avions invoqué le terme \textit{d'organisation de guérilla anarchique} pour décrire l'organisation de mise à disposition des services de piraterie et des lieux de rencontre et d'échange des communautés patrimoniales pirates. Découpons ce terme en deux éléments : \og organisation de guérilla \fg et \fg organisation horizontale \fg.

\begin{itemize}
    \item \textit{L'organisation de guérilla} : Le terme d'organisation de guérilla est une analyse produite par Martin Paul Eve dans son étude sur l'organisation minimaliste des bibliothèques pirates qu'il étudie comme Sci-hub et Libgen \footnote{Martin Paul Eve, \textit{Lessons from the Library: Extreme Minimalist Scaling at Pirate Ebook Platforms}, dans Digital Humanities Quarterly, Providence, vol n°16, n°2, 2022}. 

    \begin{quote}
         \textit{\og Ces archives, qui violent les droits d'auteur, contournent les murs payants et offrent un accès à tous les visiteurs. Bien qu'ils soient souvent du mauvais côté de la loi, les opérateurs de bibliothèques fantômes conçoivent leurs sites en termes éthiques. En définissant leur banditisme en termes de Robin des Bois, ces sites croient qu'ils volent aux riches pour donner aux pauvres de la recherche. Ce qui est important pour les thèmes de ce numéro spécial, c'est que je pense que de telles archives peuvent également être comprises en termes de minimalisme informatique. Comme je le détaillerai plus loin, divers principes de conception technique de ces archives réduisent les barrières sociales à la participation.\fg }

         Martin Paul Eve, traduit, paragraphe 7
    \end{quote}

    Cette analyse peut être mis en relation avec le \textit{Guerilla Open Access Manifesto} attribué au jeune Aaron Swartz, qui a popularisé le terme de guérilla pour décrire à la fois le mode d'organisation des communautés pirates, et la considération de lutte contre une hégémonie mercantiliste de la publication de la connaissance et de la culture.

    \item \textit{L'organisation horizontale} : Le lexique de l'horizontalité fait référence à l'organisation autogestionnaire et l'idée de libre association permettant l'émergence d'une organisation fluide permettant la survit et la réplication des lieux d'échanges des communautés pirates. L'important est la communauté plutôt que le lieu où elle se réunit, n'importe qui peut se revendiquer de la communauté du moment qu'il obéit aux règles parfois tacites, parfois écrites de cette communauté. 
\end{itemize}

L'organisation et le mode de fonctionnement des bibliothèques mis en ligne par les communautés patrimoniales pirates ont un double objectif : être à la fois simples et minimalistes tout en offrant un minimum de barrières à l'entrée et des fonctionnalités avancées et utiles aux utilisateurices. Le minimalisme provient à la fois d'un objectif de réduction des coûts de développement, de maintenance et d'hébergement des bibliothèques pirates, mais également d'offrir le moins de surface visible possible. Les bibliothèques pirates comme Sci-Hub, LibGen, Zlib se présentent souvent comme une simple barre de recherche sur un fond blanc, sans besoin d'authentification ou d'inscription utilisateur. Ce minimalisme est également un mécanisme de défense, le but d'une bibliothèque pirate est d'être le plus camouflé possible, tout en offrant la plus grande qualité et quantité de services possibles pour être partagé le plus globalement possible. Un site relativement simple sera plus facile à remettre en ligne si un domaine, du matériel, des individus sont éliminés ou emprisonnés par les autorités judiciaires. Nous reviendrons sur l'organisation des dépôts des communautés patrimoniales pirates dans le chapitre suivant, mais l'important est la différence de philosophie qui habite le design et l'utilisation des bibliothèques pirates comparée aux applications et aux services légaux. Trouver et utiliser ces outils demande des codes sociaux, des habitudes différentes, et les premiers pas dans cet univers sont souvent anxiogènes et peuvent faire tourner les talons aux individus les moins bien accompagnés.

\begin{quote}
    E - \textit{Mais plus c'est difficile d'accéder, plus je vais avoir peur d'aller fouiner, finalement. D'aller chiner des trucs, parce que je me dis que c'est caché pour une raison, je ne sais pas. Mais sinon, au-delà de la peur, c'est vraiment du soulagement, de l'euphorie ? Euphorie c'est peut-être beaucoup, j'avais besoin d'un film, je l'ai trouvé. Ça fait un peu un saut de dopamine. [...] Maintenant je n'ai plus trop peur, malgré ... Le souci des sites de streaming, c'est qu'il y a tout le temps des pubs. Des fois, j'ai un peu peur d'avoir un virus qui s'installe. Ça, c'est toujours un petit peu ... ça je n'aime pas trop.}
\end{quote}

\section{Du partage des émotions aux communautés pirates }

Dans son livre \textit{Du numérique aux émotions}, Emmanuelle Bermès adapte la typologie des émotions patrimoniales de Daniel Fabre au contexte des communautés se réunissant autour d'objets numériques. L'émotion mise en avant pour comprendre la mentalité des mouvements \og hacktivistes \fg qui mettent en avant le partage radical de l'information scientifique est la sédition. L'envie de rompre avec un système considéré comme injuste et dysfonctionnel permet de réunir des individus autour de l'idée de créer l'alternative au sein de l'infrastructure, par la mise en place d'outils, l'autogestion et une pratique révolutionnaire (ne pas jouer selon les règles du jeu) ; plutôt que d'attendre la lente évolution de la superstructure.

\subsection*{Du pirate à l'archiviste rebelle : de la consommation à une pratique politique}

Dans les entrevues avec nos enquêtés, la question du rapport politique de la pratique du piratage n'a pas été une piste ayant produit des réponses très concluantes. Nous avions déjà cité l'exemple d'Hilaire, de ses considérations éthiques et de sa découverte de la culture du libre dans le chapitre précédent. Dans le cas de Théo, la piraterie culturelle s'inscrit dans une vision globale d'un idéal d'accès à la culture ; cependant, il existe un fossé entre une vision du monde et un engagement politique personnel que Théo se refuse à franchir. De son côté, Zakary a produit des réflexions sur l'éthique de sa pratique, considérant l'organisation économique du monde des auteurs de mangas. 

\begin{quote}
    Z - \textit{Si je dois payer un élément, je dois être sûr en tant qu'acheteur, en tant que consommateur, que de mon côté je vais avoir un produit de qualité... Et un produit, du coup, assez éthique. Je dois juste être sur de la qualité, et de l'autre côté de l'artiste, il y a un traitement éthique. Il faut que l'artiste soit payé convenablement, que l'argent que je lui fournis soit bien touché, etc. Et tout simplement j'ai l'impression que moi, en tant que pirate, ce contrat a été rompu par toutes les plateformes qui proposent des alternatives légales : Netflix, Amazon, Crunchyroll, ADN, Manga Plus. Mais c'est toujours un peu critiqué dans la communauté.}
\end{quote}

Cependant, aucun de nos enquêtés ne revendique la piraterie culturelle comme la base d'une identité politique, même si certains d'entre eux se considèrent comme des personnes politisées et militantes. Dans différentes communautés patrimoniales, la sédition et la consommation pirate revêtent une couleur politique appelant à remettre en cause l'organisation capitaliste de certaines ressources. L'exemple réapparaissant régulièrement est la disparition d'un contenu culturel de la vente physique et des services de location. Certains mouvements se sont structurés autour de la lutte à offrir une vie aux objets culturels après l'exploitation commerciale, pour que la culture ne soit pas à la main des ayants droit. Par exemple, le terme d'\textit{abandonware} correspond à un logiciel qui n'est plus maintenu ni commercialisé par ses ayants droits. D'une certaine manière, l'horizon de tout objet numérique est de tomber dans l'oubli et de disparaître ou de devenir l'équivalent d'un abandonware, un objet culturel orphelin. La conscience de l'existence de la perte d'accès aux contenus culturels existe, sans réellement porter de temps. Au-delà de la question d'accès à la science et à l'éducation, à la barrière monétaire d'accès à la culture, la peur de l'oubli collectif des contenus culturels est l'une des motivations les plus profondes des communautés patrimoniales pirates.

\begin{quote}
\textit{    > Check Amazon, Netflix, Amazon, Disney, Max, Apple, etc.}

\textit{    > No where to find Yukio Mashima films
}

\textit{    > No where to find Spaceballs
}

\textit{    > No where to find Amadeus
}

\textit{    > No where to find Samurai (1954)
}

\textit{    > No where to find Jhin-Rho 
}

\textit{    They are just not "available", Why ?
}

    Anonymous n°210962625, \og My country has a Kino embargo \fg, sur 4chan, le 28/05/2025, via le post de u/Vor\_Mor, \og  No kino movie \fg, sur r/greentext, 28/05/2025.
\end{quote}

La base philosophique de la piraterie culturelle des individus et des communautés qu'ils composent est clairement d'inspiration libertaire. Cependant, il serait erroné de considérer que les forums d'échange de films sont peuplés en majorité de militants anarchistes. %A compléter


\subsection*{Les espaces numériques de discussions}

Les forums dédiés à la piraterie culturelle agissent de la même manière qu'une agora, où certains événements font office de catalyseurs amplifiant les émotions individuelles. Ces espaces de discussion et de partage deviennent le berceau d'actions concrètes. Les subreddits r/Piracy et spécifiquement r/DataHoarder sont des sources fascinantes sur les communautés patrimoniales numériques. Premièrement, ce sont des forums ouverts ; ces dernières années, une grande partie des communautés de fans, de travaux bénévoles sur des projets open source, des wikis, de communautés patrimoniales pirates a migré sur des logiciels de discussion communautaire semi-fermés comme Discord. C'est un outil fabuleux pour offrir un espace d'échange et construire une communauté ; cependant, sa nature d'espace d'échange privé a pour conséquence de supprimer du web indexable sur un moteur de recherche une grande partie des ressources et des échanges qui documentent l'existence et l'évolution de ces communautés. Deuxièmement, ces subreddits possèdent des outils de recherche intégrés comme le tri par ordre chronologique, par popularité des posts et des réponses, des controverses (les posts ayant le plus de réponses mais peu de \og upvotes \fg, l'équivalent des likes sur reddit). Ceci nous permet de suivre les discussions, les débats et les intérêts des communautés patrimoniales pirates.

\begin{quote}

     \textit{I think there will be a Jan 6 situation where this will get wiped off the internet, are there any current efforts to archive footage and images from this current ongoing event? If not I'd think that's something that should be payed attention to at the moment.} 

     u/awolfwearingabanana, \og Any attempts to archive the current LA protests? \fg, sur r/DataHoarder, le 09/06/2025
\end{quote}

Ce poste de u/awolfwearingabanana, sur l'archivage amateur de la documentation et des vidéos des émeutes ayant cours à Los Angeles en juin 2025, nous permet de plonger directement dans les discussions et les débats qui animent la communauté des \textit{archivistes rebelles} de r/DataHoarder.  Très rapidement, les premiers commentaires des utilisateurs les plus actifs font un parallèle avec les émeutes du Capitole du 6 janvier 2021, et partagent les liens torrent vers les dépôts des utilisateurs qui se sont donnés comme mission d'archiver la documentation de ces émeutes. L'expérience de la communauté des archivistes d'images et de vidéos de manifestation et les différentes méthodologies de sélection et d'extraction de vidéos sont discutées sous ce poste. D'autres creusent la distinction entre l'amassement de données, l'archivage et la mise à disposition du public dont ces données auraient besoin. L'initiative ne fait cependant pas l'unanimité, et une partie des membres du forum critique la propension à une forme de paranoïa archivistique, un complotisme de la mémoire. 

\begin{quote}

\textit{My suggestion would be to download content you find. Assume others are not saving anything, then share your collection.
There's a 1TB torrent of Jan6 content which I believe was a compilation of files different people downloaded and then shared here.
(edit: I'm not from the US or saving any content. Just suggesting what you can do and pointing out what others done before.)}

u/P03tt, en réponse à u/awolfwearingabanana sur r/DataHoarder.\\

\textit{Archival is the first step, and what a lot of people do, but I would argue it's just the first step. If you're looking at preserving data for the public good, you should look at making those archives available, then curate and index them. If you do create an archival site, feel free to let us know and we can add it to the https://datahoarding.org/ index. We already have a lot of Jan6 archives on there. }

u/shimoheihei2, en réponse à u/awolfwearingabanana sur r/DataHoarder.\\

\textit{Jan 6 situation where this will get whipped off the internet
Except for the fact that this never happened? Everyone knows what Jan 6 is, all the videos and photos are online and easily accessible.
I swear redditors live in this weird meta-world where they believe reddit is simultaneously some kind of safe harbor for talking about events the rest of the world would get "removed" for. If "they" (whoever you believe wants to "whip" things of the internet) is able to actually do so, why do you believe this post is up? If what you believe was real you'd be taken into a black site and "questioned" while your post got removed the second it hit a server. payed .Where's the bot? Oh wait, no, I get it. You're on a boat stuck with very slow satellite internet and believe everything is being censored because it won't load on your shitty connection?}

u/IronCraftMan, en réponse à u/awolfwearingabanana sur r/DataHoarder.\\

\end{quote}

\section{Conclusion}

L'horizon mental et émotionnel des individus qui piratent peut être étudié sous le prisme des théories de la sociologie de la déviance. La piraterie culturelle, avec les définitions de la sociologie, peut être considérée comme un fait normal \footnote{Albert Ogien, \textit{Le normal et le pathologique}, dans \og Sociologie de la déviance \fg, p. 29-38, Presses Universitaires de France, 2012}, que l'on retrouve uniformément dans des contextes séparés. L'élément de réflexion qui rend l'étude de la piraterie culturelle très intéressant est la jonction entre un acte de déviance (je refuse de payer pour un objet numérique payant) et ses conséquences (peur, clandestinité, prise de risque), une revendication politique et culturelle (conservation, éducation, rapport aux institutions et entreprises ayant les droits d'exploitation), et un aspect communautaire spécifique à la dimension patrimoniale (forums, organisation de partage et de collecte, bibliothèques pirates). La massification de la capacité de copie des individus entre en contradiction avec le système de distribution culturelle et les règles qui l'entourent. Le sentiment de dépossession face au patrimoine numérique pousse vers la sédition et le piratage comme catalyseur d'alternatives. Cet espace clandestin demande un apprentissage social et culturel pour pouvoir être navigué, mais il renforce les liens communautaires et permet aux individus de se sentir maîtres de leur consommation culturelle en forgeant de nouvelles compétences au cours de leurs carrières pirates. L'apprentissage d'une culture de la piraterie culturelle permet aux individus de se regrouper en communautés, avec comme moteur les émotions générées par le rapport de la société aux objets culturels immatériels. Ces considérations pour des objets culturels rentrent dans le domaine du politique, même si les individus ne sont pas forcément des militants politiques. 


\chapter{Bibliotheca mundi, Hic sunt dracones}

\begin{quote}
    \textit{L'univers (que d'autres appellent la Bibliothèque) se compose d'un nombre indéfini, et peut-être infini, de galeries hexagonales, la distribution des galeries est invariable. Vingt longues étagères, à raison de cinq par côté, couvrent tous les murs moins deux. Chaque étagère comprend trente-deux livres, tous du même format. La Bibliothèque est totale, et ses étagères consignent toutes les combinaisons possibles des vingt et quelques symboles orthographiques (nombre, quoique très vaste, non infini), c'est-à-dire tout ce qu'il est possible d'exprimer, dans toutes les langues.}

    Jorge Luis Borges, \textit{La Bibliothèque de Babel}, dans \og Fictions \fg, 1944
\end{quote}

Dans nos deux chapitres précédents, nous nous sommes d'abord interrogés sur le début des carrières pirates, puis sur le lien entre le vécu des individus et la création de communautés patrimoniales pirates. Dans ce troisième chapitre, nous édifierons une typologie d'étude des actions de ces communautés, en étudiant les dépôts \footnote{Pour joindre les conceptes d'archives, de bibliothèques et de tous les espaces communautaires ne tombant pas exactement dans ces définition, Abigail de Kosnik utilise le terme de répertoire. En réalité, son utilisation du terme répertoire est plus vaste que ce que étudions ici (il s'agit d'un répertoire culturel), mais je pense que l'utilisation du mot dépôt permet de bien comprendre la nature de ces tiers-lieux clandestins.} que créent et utilisent ces communautés. La manière dont est conçue une plateforme de communication et ses fonctionnalités peut grandement influencer la manière dont les individus échangent l'information. Par exemple, dans les archives des skyblogs de la BnF, nous retrouvons plusieurs blogs dont le but est d’offrir un accès à des fichiers et des logiciels, mais il est difficile d'observer une dynamique communautaire concernant la piraterie culturelle, ses conséquences politiques et ses mises en place infrastructurelles. Sur les plateformes de blogging et de microblogging comme Skyblog et Twitter/X, les dynamiques communautaires sont atténuées par le fonctionnement individualiste de partage des plateformes. Pour réussir à être visible sur Twitter, Facebook, Skyblog, ou la plupart des réseaux sociaux modernes, il faut avoir une base de followers, des gens qui se sentent touchés et investis par les messages, les blagues et les informations qu'un.e utilisateurice particulier fait circuler. Reddit et Discord sont quant à eux des réseaux purement communautaires, où les discours d'une personne peuvent toucher toute une communauté sans avoir besoin de construire une notoriété préalable. Nous étudierons ici les différents modes de partage, d'organisation et de maintien des dépôts pirates de différentes communautés.

\section{Naviguer sur les sept mers : une visite expresse des dépôts pirates personnels.}

Pourquoi ne pas utiliser le terme de bibliothèque pirate comme le fait l'article \textit{Lessons from the Library: Extreme Minimalist Scaling at Pirate Ebook Platforms} de Martin Paul Eve ? Ou pourquoi ne pas utiliser le terme de \textit{Rogue Archives} comme l'utilise Abigail de Kosnik dans son livre éponyme ? La raison de l'utilisation d'un nouveau terme à définir apparaît à la nature d'un constat, les archives des un.es peuvent devenir les bibliothèques et les répertoires des autres. Nous nous concentrerons ici sur la cartographie du partage interpersonnel, en rencontrant des individus qui partagent leurs collections de biens culturels au gré de leurs envies. S'ils peuvent faire partie d'une communauté patrimoniale, la manière dont ils mettent en place leurs systèmes de partage de l'information reste un élément personnel. 

\subsection*{Le dépôt personnel}

Le dépôt pirate personnel est, comme son nom l'indique, créé, alimenté et maintenu en ligne par une seule personne. Cette personne peut choisir de maintenir ce dépôt pour son profit propre, ou de le partager ; comme votre bibliothèque de salon, le dépôt personnel est défini par l'unicité de son origine et non par le nombre de personnes qui y ont accès. Ces dépôts peuvent être un simple dossier où l'individu stocke ses livres ou ses jeux piratés, ou des dispositifs de stockage beaucoup plus complexes et onéreux.

\begin{quote}
    T - \textit{[Avec mon frère] au début, on s'est échangé des clés USB ou des disques durs pour se partager des films qu'on avait envie de voir du registre de l'autre. Mais c'était devenu assez fastidieux, et surtout on c'est éloigné géographiquement pour avoir chacun sa vie pro dans différentes villes. [...] Donc il s'est posé la question de pouvoir avoir accès au registre de l'autre et de tout mettre en commun. On s'est renseigné, c'était une possibilité pour nous d'avoir notre propre serveur en fait. Récemment on a lancé ce projet à fond. On est encore en train de faire des testes et tout, mais on a ce qu'on appelle un NAS \footnote{network attached storage}. Ce serveur c'est mon ancien pc qui a été dépiécé et un peu agencé pour une nouvelle utilisation. On a mis de quoi rajouter plein de disques durs avec des énormes capacités, donc en l’occurrence on a en totalité ... La capacité totale, on a 4 disques de 8 Tera, ça fait 32 \footnote{32 000 gigaoctets}. [...] Pour nous ça fait un peu court, parce qu'on savait que nos deux bibliothèques, nos deux registres, on va dire cumulés, on était déjà à plus de 4 ou 5 Tera, et là actuellement, le dernier chiffre aujourd'hui on est presque à 7 tera au total.}
\end{quote}

Lors d'évènements particuliers, le dépôt personnel peut être ouvert par son propriétaire pour en faire profiter son cercle, ou toute une communauté. Dans le cadre de Théo et de son frère, le dépôt créé contient tellement de films et la logistique matérielle et logicielle le permettant, qu'ils décident de l'ouvrir aux membres de la famille et à leurs ami.es.

\begin{quote}
    T - \textit{Nous, TrueNAS, c'est un choix de facilité car c'est gratuit, et au dessus on a mis Jellyfin qui nous sert de lecteur média depuis ce système-là. [...] On a regardé la carte graphique de mon ancien pc, c'est une GTX 950. Donc ça date, ce n'est pas vraiment très puissant, mais il se trouve que c'est largement suffisant pour faire du décodage de 5 lecteurs en même temps. [...] On ouvrirait à des amis et on leur ferait payer à l'année une petite somme, quelque chose comme 10 euros, simplement pour compenser la facture d'électricité d'avoir un pc qui tourne tout le temps. La deuxième motivation c'est pour nos parents, nos parentes étaient très dépendants, ils avaient opté pour la facilité, c'est que te propose par exemple Netflix. Pour eux c'est la même chose, la différence, c'est qu'ils ne paieront plus Netflix, et nous derrière on aura plus de contrôle sur ce qui se passe sur le catalogue pour éviter que des trucs disparaissent du jour au lendemain.}
\end{quote}

D'autres exemples moins technophiles peuvent être intéressants à étudier. Nous avions évoqué les vagues de partage de films de David Lynch. À l'annonce de la mort du réalisateur le 16 janvier 2025, de très nombreuses personnes se sont mises à télécharger, riper et partager les œuvres du réalisateur. En ouvrant son dépôt personnel en le transférant sur son Google Drive ou un MegaUpload, les individus prenaient part au deuil international de la communauté des fans. Le piratage et son partage revêtent ici à la fois une dimension mémorielle (les œuvres ne disparaissent pas si les fans en créent des copies) et communautaire (je peux contribuer à rendre hommage à David Lynch en propageant son art).

\begin{quote}
\textit{    coucou j’ai téléchargé toutes les œuvres de Lynch (films et série) et je les ai mis dans un mega pour que ce soit plus facile pour toutes les personnes qui veulent (re)découvrir ce qu’il a fait :) et n’oubliez pas, si vous commencez twin peaks l’ordre c’est : saison 1 —> saison 2 —> Fire walk with me —> saison 3
}

    @cycylcds, sur twitter/X, le 18/01/2025
\end{quote}

\subsection*{Les individus \og passerelles \fg et la culture archivistique amatrice}

Il existe un entre-deux entre le dépôt personnel et le dépôt communautaire. Il s'agit de dépôts dont l'administrateur est une personne particulière, mais dont le but de création est d'être partagé à une échelle communautaire. Ces dépôts sont parfois éditorialisés, par des individus qui sont reconnus dans leurs communautés comme étant des ressources précieuses de savoir et de connaissance. Ce genre d'individu que nous pouvons surnommer des \og passerelles \fg sont relativement courant dans les communautés politisées comme les milieux militants et LGBT. Ce partage passe par des infrastructures et une organisation plus pérennes que les simples dépôts personnels. Nous pouvons trouver dans cette catégorie des exemples divers, allant de la subversion de systèmes électroniques pour la réparation d'objets, à la mouvance du \textit{bio-hacking} et de l'Anarcho-transhumanisme transgenre refusant d'être dépendant.es de l'ordre médical et de l'industrie pharmaceutique pour leurs transitions de genre. 

\begin{quote}
    \textit{They're both on my website now!! (Une image représentant deux fascicules) Transfem DIY HRT, Transmasc DIY HRT, Nobody can stop you}

    @desert\_varnish, sur twitter/X sur 02/02/2025
\end{quote}

La principale communauté de personnes \og passerelles \fg que nous avons pu étudier se situe sur le forum subreddit r/DataHoarder. Nous l'avons déjà citée dans le chapitre précédent sans vraiment expliquer en profondeur par quoi cette communauté se passionne ainsi que comment interagissent ces membres. Le terme de datahoarder provient du mélange du terme data (données) et hoarder (ceux qui stockent). Cette communauté est composée de personnes fascinées par le stockage, l'archivage et la mise à disposition de jeux de données. Les datahoarders mettent en avant les qualités de pérennisation et de partage de leurs bases de données. Cependant, le choix des données qu'ils stockent et la mise à disposition est un choix personnel, qui suit souvent les intérêts et passions du hoarder. r/Datahoarder est donc un rassemblement de personnes stockant et partageant des collections de données rares, mais le fonctionnement interpersonnel (demander à un hoarder s'il compte rendre ses données publiques) fait que nous sommes face à une multitude de bibliothèques personnelles plutôt qu'à une communauté organisée autour d'un outil.

Il est passionnant de se plonger dans les interactions d'une communauté si riche qui se base sur l'interaction sur un forum ouvert pour communiquer. Par exemple, le forum r/DataHoarder permet aux membres de créer une étiquette qui s'affiche après le pseudonyme ; ils s'en saisissent afin d'afficher les différentes tailles de leurs collections et infrastructures. Il se dégage ainsi une forme de hiérarchie immanente au volume qu'une personne stocke et à la logistique et l'argent qu'iel met en place pour mener à bien sa mission de hoarder. Certaines étiquettes laissent deviner une forme de syndrome de Diogène numérique, avec des personnes indiquant avoir mis en place des serveurs de stockage de plusieurs centaines à plusieurs milliers de téraoctets de stockage. Si beaucoup de ces utilisateurs sont des ardents défenseurs du stockage et du piratage de contenus culturels, beaucoup ne se considèrent pas comme des pirates. La citation ci-dessous est particulièrement intéressante car u/Blue-Thunder utilise la terminologie académique du livre \textit{Rogue Archives} d'Abigail de Kosnik, en modifiant la définition qu'iel applique derrière \footnote{Agigail de Kosnik utilise ce terme de \textit{Rogue Archivist} pour parle principalement des personnes archivants et partageants des contenus créer par des communautés de fan (fanfictions, fanarts), ici u/Blue-Thunder archive des fichiers qui sont des dossiers sous droit, comme des films, des livres, des vidéos, ce qui est très différent du première exemple. }. 

Un excellent exemple trouvé par Marina Hervieu avec les outils de recherche de la Bnf est le skyblog \textit{brotjackta} de l'utilisateur Ruinessi. À première vue, ce skyblog est rempli de posts vides qui affichent la phrase suivante \textit{Sorry, the page you requested has been moved.}. Les titres des posts sont pourtant très intéressants, il s'agit de tutoriels anglophones, pour découvrir plusieurs outils de rapatriement de données torrents ou de téléchargements. En réalité, le blog de Ruinessi est écrit en blanc sur fond blanc. Avec cette astuce, il espère éviter d'attirer l'attention des personnes n'ayant un minimum de connaissances en navigation web et informatique. Nous avons trouvé d'autres exemples d'individus assurant un rôle de passerelle dans les archives des skyblogs. Les posts et les blogs ont la particularité d'être très simplement personnalisables, avec notamment du HTML, et de pouvoir sauvegarder une petite quantité de données extérieures comme un fichier mp3 ou quelques images. Certains utilisateurs ont profité de ces possibilités, couplées à la capacité d'offrir des hyperliens vers des sites extérieurs pour offrir des mini-bibliothèques pirates.

Voici quelques exemples de blogs de bibliothèques pirates de la plateforme Skyblog.
\begin{itemize}
    \item \textit{FilmstreamingVF}, créé en 2016 par l'utilisateur telechargerfilm, probablement lié à un site de téléchargement filmvfr.com. Ce blog utilise des petites fiches d'informations techniques, ainsi que des critiques de différents films de l'actualité cinématographique pour amener ses utilisateurs à passer sur son site de streaming pirate.
    \item \textit{téléchargergratuitement}, créé en 2018 par un utilisateur s'inspirant du fonctionnement de 01.net,  qui est un peu particulier ; en effet, ce site propose de télécharger des applications de manière légale. Son modèle de fonctionnement s'inspire des applications d'installation comme les différents gestionnaires de paquets sous Linux. Chaque article de blog s'intitule selon une application, avec une description de son utilité et de son usage. 
    \item \textit{La biblio de Tetel} et \textit{Le monde de Tetel}, deux blogs liés et créés par l'utilisateur Kurisu/Tetelle dont la première capture sur la Wayback Machine remonte à 2013. Cette personne reprend les codes de l'article de blog en donnant un avis critique sur plusieurs œuvres, en ajoutant des liens de téléchargement dans le corps du texte ou dans les commentaires. 
    \item \textit{Mangas-a-telecharger}, créé en 2008 par un utilisateur dont je n'ai pas trouvé le pseudonyme. Ce blog est intéressant car il s'agit du blog de téléchargement pirate le plus assumé de nos trouvailles. Après une petite introduction, l'auteurice du blog affiche une liste de liens pour avoir accès à ses mangas favoris. Quelques posts de blog demandent de l'aide à ses abonnés pour trouver les mangas sur lesquels iel n'arrive pas à mettre la main. 
\end{itemize}

Nous pouvons également trouver quelques initiatives solitaires comme le blog \textit{Anti-téléchargement} \footnote{https://web.archive.org/web/20130517194237/http://anti-telechargement.skyrock.com/} essayant de prévenir et de sensibiliser aux dangers du piratage, ou le blog \textit{Unis-Contre-Le-Plagiat} \footnote{https://web.archive.org/web/20100215052912/http://unis-contre-le-plagiat.skyrock.com/} sur la question du plagiat de posts entre utilisateurs de Skybklog. \\

La plupart de ces blogs prennent souvent la même forme. Un post de blog avec comme titre le nom de l'œuvre (souvent des mangas ou des films) avec ensuite un résumé et une critique de l'œuvre, puis une note et enfin un lien de rapatriement du fichier.\\

%ajouter un élément qui introduit cette citation au mon dieu 

Sur r/DataHoarder, les habitués du forum se définissent eux-mêmes comme des archivistes rebelles, revendiquant un fonctionnement clandestin au nom d'une lutte pour le patrimoine numérique. 

\begin{quote}
    \textit{I'm sorry but I think this terminolgy is outdated. I prefer to be called a Rogue Archivist. I'm not on the high seas pludering loot from corporations, I am protecting media from vanishing forever, against the will of the respective copyright owners who would love to see certain things just vanish, or who are so incompentent that they do not have (redundant) backups.}

    u/Blue-Thunder, en réponse au poste de u/pilimi\_anna \og  How to become a pirate archivist \fg sur r/datahoarder, 2022
\end{quote}

Cependant, r/DataHoarder n'est pas dédié au stockage de données dans un but précis, seulement au stockage et à la mise à disposition de données de masse, et ceci peut amener à des discussions de questions éthiques. Si le stockage de contenus culturels piratés n'est pas considéré comme un problème ou un dilemme, le stockage et la mise à disposition de données personnelles l'est beaucoup plus. Par exemple, plusieurs projets de membres se concentrent sur le scrapping, le stockage et la diffusion de données pornographiques. Le \textit{Petabyte Porn Project} \footnote{u/-Archivist, \textit{The Petabyte Porn Problem: Public webcam social interaction archiving. [Chaturbate, MyFreeCams, Streamate]}, sur r/DataHoarder, 2017} se présente sous l'angle de l'étude et de l'archivage des \og interactions sociales \fg. Pour ce faire, les individus qui se joignent au projet via une feuille de contact peuvent archiver dans une base de données commune les fichiers vidéo de plusieurs sites de pornographie en direct \footnote{\og Le phénomène connu sous le nom de « caming », qui consiste pour des « camgirls » ou « camboys » à proposer, moyennant rémunération, une diffusion d'images ou de vidéos à contenu sexuel, le client pouvant donner à distance des instructions spécifiques sur la nature du comportement ou de l'acte sexuel à accomplir, n'entre pas dans le cadre de la définition de la prostitution qui implique un contact physique onéreux avec le client pour la satisfaction des besoins sexuels de celui-ci.\fg, citation d'un acte de la chambre criminelle de la cour de cassation, n° 21-82.283 B, 18 mai 2022} (Chaturbate, MyFreeCams, Streamate). Plusieurs utilisateurs vantent les utilisations possibles d'une telle collection : étude sociologique sur les interactions entre regardants et regardé·es \footnote{u/t0asterB0y \og This could be an invaluable resource for social scientists and behavioral psychologists. \fg, 2017}, tentative d'application de reconnaissance faciale pour retrouver des victimes de traite humaine \footnote{u/deleted,  \og Have you thought about doing facial recognition and matching up to sex trafficking photos, missing persons photos, etc? \fg, 2017}. Cependant, r/DataHoarder est une communauté essentiellement composée d'hommes, et les réponses à ce post de blog sont remplies de personnes demandant à avoir accès à la base de données pour chercher leurs stars préférées, en particulier celles qui ont arrêté de gagner leur vie avec la pornographie et ont supprimé leurs vidéos en espérant ne plus avoir à assumer une notoriété publique sur le web. D'autres éléments problématiques sont mis en avant pour les utilisateurs les plus sceptiques : la détention et le stockage de fichiers pornographiques de personnes mineures sont strictement interdits dans la grande majorité des pays, même à des fins documentaires et de tentative d'aider les autorités compétentes. 

\section{La carte et le coffre au trésor : les dépôts communautaires pirates}

Dans notre typologie des dépôts pirates, toutes les initiatives de stockage et de partage que nous avons étudiées ont comme points communs d'être des initiatives individuelles. Ces dépôts peuvent être personnels ou, au contraire, s'adresser à une communauté, mais les créateurs et les mainteneurs de ces projets étaient toujours des personnes seules, ou des petits groupes de personnes liés par des relations interpersonnelles. Cependant, ce n'est pas parce qu'un dépôt pirate est un dépôt personnel qu'il ne peut pas s'inscrire dans une communauté patrimoniale pirate. Ce qui différencie les dépôts personnels et les dépôts communautaires est à la fois le mode de fonctionnement, l'ampleur des projets et la cible de la diffusion des projets. 

\subsection*{Les dépôts communautaires pirates}

Les exemples les plus connus de dépôts communautaires sont les bibliothèques, académiques ou pas, pirates. Que ce soit les plus connues : Sci-Hub, LibGen et Zlibrary, le sujet de ces bibliothèques en ligne a été l'épicentre de nombreux débats, d'articles médiatiques et de quelques études. Les bibliothèques pirates comme LibGen sont des dépôts communautaires pour plusieurs raisons : 

\begin{itemize}
    \item Ces dépôts pirates sont trop massifs pour être gérées par un seul utilisateur, plusieurs équipes entières travaillent sur ces projets afin d'assurer la logistique, de maintenir le fonctionnement des sites et des serveurs, mais également de développer de nouvelles fonctionnalités pour les utilisateurs. Certaines de ces projets sont également open-sources, ce qui agrandit encore le cercle de personnes impliquées dans le fonctionnement. Prenons un exemple concernant le jeu-vidéo : le projet XenomRecomp du développeur Hedge-dev \footnote{https://github.com/hedge-dev/XenonRecomp} permet de recompiler des jeux Xbox 360 en des exécutables natifs pour ordinateurs Windows. Copier et recompiler du code est évidemment une pratique pirate, mais cet outil est aussi extrèmement efficace pour sauvegarder et faciliter l'utillisation de toute la librairie de la console. Les personne ayant des connaissances aboutis en développement, et de l'expérience avec la compilation ou l'émulation de logiciels peuvent aider le développeur original à fiabiliser et développer son outil grâce au système de tickets nommé \textit{Issues} intégré au site Github.

    \item Tous les dépôts dont nous avons étudiés dans cette partie n'ont pas été créé par et pour une seule communauté. Ils ont pour but d'être partager et utilisé par quiconque qui peut trouver une utilité dans leurs outils et les documents que les dépôts communautaires mettent à disposition. Plusieurs de nos exemples cités ont pour but de mettre à disposition le plus d'imprimés possibles et ce à l'échelle internationale. 
\end{itemize}

Zlibrary, comme SciHub, LibGen ou Anna's Archive, sont l'architecture de fonctionnement sur laquelle nous construisons le concept de \textit{l’organisation anarchique de guérilla} que nous avons évoqué dans notre introduction et notre deuxième chapitre. Ce sont des sites qu'il peut être difficile de trouver quand on ne sait pas où chercher. La plupart de ces exemples ne sont pas directement disponibles en France dû à des décisions de justice imposant aux principaux fournisseurs d'accès à Internet de les effacer de leurs DNS \footnote{
Domain Name Systemn ou Système hiérarchique de nommage des ressources d'un réseau : un des protocoles de base de l'Internet, sans inscription d'un site à un dns, vous ne pouvez pas vous y connecter.
}. Et certaines des URL de ces sites sont régulièrement fermées, affichant alors une image bien connue des habitués, une page bleue avec le logo du FBI américain ainsi qu'une explication de la saisie du domaine et du caractère illégal du site qui y était hébergé. Il faut alors pour l'utilisateur trouver des alternatives, les quatre principales bibliothèques ont ce que l'on nomme des miroirs, des sites jumeaux, ou des messageries automatisées sur le service Telegram, permettant de faire des demandes de retrait de fichiers sans passer par l'interface utilisateur du site. Ces solutions de repli peuvent prolonger la vie du site pendant quelques temps avant que les autorités ne remontent aux serveurs. Quand un service passe complètement hors ligne, nous pouvons observer des messages de déploration et de recherche d'alternatives sur les principaux forums communautaires (reddit) et les sites de microblogging (twitter/x).

\begin{quote}
    \textit{I miss Zlib everyday </3}

    @BLincoln56151, sur Twitter/X, le 20 avril 2025
\end{quote}

Cependant, les équipes travaillant sur ces bibliothèques d'Alexandrie contemporaine sont si vastes et organisées qu'une nouvelle version du site ne tarde pas à réapparaître en ligne après quelques semaines ou quelques mois. Cependant, il est difficile pour les équipes officielles de ne pas concurrencer des groupes malveillants de phishing, cherchant à récupérer des adresses e-mails et des mots de passe en créant des faux clones incomplets et en les faisant passer pour le site officiel.

\subsection*{Bibliothèques pirates et archives rebelles : étude de cas avec \textit{Zlibrary} et \textit{The Eye}}

\subsubsection*{\textit{Zlibrary}}
Prenons un exemple : la plateforme Zlibrary est une bibliothèque clandestine mettant à disposition des milliers de livres et d'articles sur tous les sujets. Zlibrary correspond parfaitement aux mécanismes décrits dans le chapitre précédent de simplicité apparente, avec un minimalisme de surface qui peut créer une impression d’inhospitalité à première vue. La première page du site se compose du logo de la Zlibrary, suivi du slogan \og \textit{ Your gateway to knowledge and culture. Accessible for everyone. } \fg et d'une simple barre de recherche. Mais si l'on passe cette première introduction minimaliste, Zlibrary offre énormément de fonctionnalités bien plus avancées qu'un simple catalogue de livres à télécharger. 

Premièrement, le site propose un système de compte permettant aux utilisateurices de garder une trace des livres qui les intéressent, de créer des listes de lecture, de se rappeler de ce qu'iels ont téléchargé, d'être informé·es de leur nombre de téléchargements restants, d'avoir accès à des recommandations de lecture personnalisées selon leurs intérêts. Zlibrary \footnote{https://z-lib.gl/, mais vous n'avez pas le droit d'y aller, ce serait illégal} est un outil avec une présentation minimaliste, mais qui essaie de fournir à ses utilisateurices une expérience lisse, complète, parfois plus adaptée à certains publics. Par exemple, pour le public estudiantin et les professionnels qui ont besoin d'un accès à la littérature académique, Zlibrary est un outil beaucoup plus adapté que n'importe quelle plateforme de vente d'ebook. On y trouve à la fois de grandes monographies, différents éditeurs, des revues, des recommandations.

Ensuite, on y découvre l'outil des listes de lecture créées par d'autres utilisateurices. Ces listes de lecture pourraient être un objet d'étude à part entière ; par défaut, les comptes Zlibrary sont publics et tous les utilisateurices peuvent voir ce que les autres utilisateurices regardent, classent et téléchargent. Imaginons un utilisateur type de Zlib, un étudiant travaillant sur le piratage comme vecteur d'une culture archivistique populaire. Pour garder une trace des différents livres et papiers qu'il explore sur ce sujet, il crée plusieurs listes de lecture, la première se nomme \textit{Piracy}, la seconde \textit{Digital Curation}, la troisième \textit{Digital Humanities}. Il utilise alors Zlibrary comme une bibliothèque, afin de mener à bien une recherche bibliographique. Une fois son travail de recherche bibliographique terminé, ses documents téléchargés, ses listes de lecture publiques perdurent, elles deviennent des archives numériques, des traces de son sérieux académique et de son érudition. Cependant, le caractère public et indexé dans le moteur de recherche des listes permet de leur donner une seconde vie. Plusieurs autres utilisateurices vont également faire des recherches sur ces sujets pour leurs études ou par curiosité, et alors les archives des uns deviennent les bibliothèques des autres. Au total, archives de lecture de notre étudiant exemple ont été consultées 586 fois. Si Zlibrary offre des outils de classification et de recherche avancée, les listes de lecture sont une fonctionnalité précieuse qui permet de se reposer sur le travail bibliographique d'autres personnes pour accéder à des documents moins populaires, moins cités, ou qui ne sont pas considérés comme pertinents par l'algorithme de recherche.\\

\subsubsection*{\textit{Archive Team \& le projet The Eye}}

Archive Team \footnote{https://wiki.archiveteam.org} est une communauté en ligne ressemblant beaucoup dans ses missions et son fonctionnement à r/DataHoarder. Le site d'Archive Team se présente comme une équipe de \og d'archivistes rebelles \fg, de développeurs et d'écrivains qui se mobilisent pour préserver l'histoire et l'héritage numérique depuis 2009. Cette communauté utilise et reformule la définition de \textit{d'archivistes rebelles} définie par Abigail de Kosnik. Comme r/DataHoarder, les projets d'Archive Team se sont construits sur la réflexion de l'enjeu culturel et politique de la conservation de certaines données par des groupes qui ne sont ni des institutions, ni des entreprises. Le projet \textit{The Eye} mis en place par l'Archive Team est cependant très différent d'un dépôt pirate en centralisant les actions de ses membres en un seul collectif. Si le projet d'Archive Team se présente comme étant un groupe respectant les règles établies de la propriété intellectuelle, l'un des projets issus de cette communauté, \textit{The Eye} s'est fait connaître pour son archivage et sa mise à disposition de dépôts de torrents pour l'entièreté de Library Genesis (LibGen) en 2019 et nous pouvons également retrouver dans l'arborescence \textit{../files/comics/..} des archives complètes de mangas sous droits, mais également des manuels de jeu de rôle, des articles scientifiques pour de l'entraînement DIY de modèles d'intelligence artificielle, et des collections comme \textit{The Telegram [Anti]Social Research Archives}\footnote{https://the-eye.eu/tasra/} avec plus de 200 groupes de discussions publics sur la plateforme Télégram sélectionnés pour la virulence de leurs membres avec presque 4 To de données réunies. Le site reprend énormément du fonctionnement et de l'identité graphique des dépôts pirates qui revendiquent leur appartenance à la mouvance des dépôts pirates. Le simple nom de la communauté peut sembler inquiétant : L'œil, avec un logo représentant un œil ouvert, omniscient, au centre d'un triangle équilatéral où chaque segment est marqué de trois marques. La page d'accueil est composée d'un fond noir, avec en transparence des paires d'yeux dessinés avec une direction artistique ésotérique et les quelques lignes de texte écrites ne permettent pas à l'utilisateur néophyte de continuer sur le site pour en apprendre plus sur l'organisation, le fonctionnement et les projets en cours de The Eye et de l'Archive Team. Cette présentation minimaliste récurrente, couplée à des directions artistiques souvent ésotériques, monochromes, chaotiques, encourage une expérience utilisateur de curiosité et de jeu avec un interdit réel ou fantasmé \footnote{Martin Paul Eve, \textit{warez: the infrastructure and aesthetics of piracy}, édition puctum books, 2021}. L'esthétique et le politique se mélangent ici, pouvant notamment nourrir la manifestation d'une représentation du monde que nous avions définie comme un zèle ou un complotisme archivistique dans le chapitre précédent. 

\section{Conclusion}

Dans ce chapitre, nous avons produit un panorama des différents types de dépôts pirates, en discriminant entre les différentes architectures et organisations des individus qui les créent et les maintiennent. Nous avons appliqué les grilles de lectures et d'analyses de Martin Paul Eve qui mélangent les modalités d'organisation de ces dépôts avec une suite de codes culturels, une esthétique du piratage. Ceci nous a permis de comprendre la traduction du principe d'organisation de guérilla et les influences sous-jacentes qui en font un référentiel d'organisation autant technique que politique. Ce référentiel est composé d'un discours composé de plusieurs sphères d'influences. Le cœur de ce discours est en grande majorité alimenté par la culture du libre et ses théoriciens/ingénieurs principaux comme Richard Stallman, Aaron Swartz ou Linus Torvalds. D'autres éléments peuvent être empruntés à la culture militante marxiste, anarchiste et libertaire. Comme toutes les cultures qui remettent en cause le statu quo, la culture technique et politique des individus baignant dans le monde des dépôts pirates peut amener à des tensions internes. Certains individus désirent conserver tout ce qui est à la portée de leurs capacités de stockage, y compris des données personnelles sensibles, ou en entrant dans une forme de \og complotisme archivistique \fg. Le fonctionnement, l'organisation et le rôle de ces dépôts pirates peuvent être interprétés comme le symptôme d'une société qui cherche à se défendre face à l'organisation du travail et du rôle de la marchandise dans son fonctionnement. Nous pouvons citer ici les grandes lignes des théories de la littérature marxiste sur le sujet de la culture, sa place dans la superstructure et son lien avec l'infrastructure. À l'époque de la reproductibilité technique parfaite de la grande majorité des objets culturels, ces dépôts pirates, les individus et les communautés qui les maintiennent, jouent un rôle important dans la sauvegarde et la diffusion de savoir et portent en eux la graine d'un possible d'une distribution plus juste de la culture et d'une collectivisation de l'information. 


%En conclusion, les dépôts pirates, qu'ils soient personnels ou communautaires, jouent un rôle essentiel dans la sauvegarde et la diffusion du savoir à l'ère numérique. Les dépôts communautaires, tels que LibGen, Sci-Hub et Zlibrary, démontrent une capacité impressionnante à s'organiser et à prospérer malgré les défis juridiques et techniques. Ces bibliothèques pirates représentent non seulement des efforts collectifs importants pour contourner les limitations imposées par les systèmes traditionnels d'accès à l'information, mais elles suscitent également des débats persistants sur la légitimité et l'éthique de leur existence. 

%En conclusion, la plateforme Zlibrary ainsi que des initiatives comme Archive Team et The Eye illustrent une approche unique et audacieuse de l'archivage numérique et de la diffusion du savoir. Zlibrary, avec son minimalisme apparent et ses fonctionnalités étendues, se présente comme une ressource précieuse pour ceux qui recherchent un accès universel à la connaissance. Simultanément, Archive Team et The Eye s'engagent dans la préservation et le partage des informations, tout en remettant en question les paradigmes traditionnels de la propriété intellectuelle et de l'archivage. Ensemble, ces plateformes témoignent d'une mouvance qui mêle esthétique, politique et une certaine subversion pour rendre accessible une vaste gamme de contenus. Cependant, cette accessibilité soulève des questions éthiques et légales qui méritent d'être examinées de près.


%%III/ Dépôts numériques, hybride entre un lieu de
%conservation et de partage.
%• Le dépôt pirate, entre archive et bibliothèques
%– La fonction de conservation des dépôts numériques.
%– Le partage et la diffusion des ressources piratées.
%• Les communautés autour des dépôts numériques
%– Les contributeurs et les utilisateurs des dépôts.
%¨ Les dynamiques sociales et les interactions au sein des communautés.
\chapter{Conclusion générale}

\begin{quote}
    \textit{On peut, bien entendu, comprendre de la même façon le développement de l'art, et de la science, de la religion et de toutes formes de culture, élevées ou inférieures. En tirant une leçon des efforts des uns et des autres et en appréciant leurs nombreuses contributions, les êtres humains construisent progressivement des systèmes de connaissances et de croyances ; ils élaborent des techniques reconnues et des styles élaborés de sensibilité et d'expression. Dans ces exemples, le but commun est souvent profond et complexe, étant défini par la tradition correspondante, artistique, scientifique ou religieuse et, pour le comprendre, il faut souvent des années d'étude et de discipline. L'essentiel est de partager  une fin et d'être d'accord sur les moyens de l'atteindre afin de permettre la reconnaissance publique des réalisations de chacun. Quand cette fin est atteinte, tous trouvent de la satisfaction dans la même chose ; et ce fait ainsi que la complémentarité du bien des individus renforcent les liens de la communauté.}

    John Rawls, \textit{Théorie de la justice}, chapitre 79 \og L'idée d'union sociale \fg, page 569, éditions du Point, 1971.
    
\end{quote}

Voici venu l'instant de clore ce travail sur l'univers foisonnant de la culture du piratage culturel, des individus qui vivent cette culture au jour le jour, des communautés qu'ils forment, et des actions que ces communautés entreprennent pour faire vivre l'idée que tout ce qui a été digne d'être créé mérite d'être conservé et accessible à tous. 

Dans notre introduction, nous avons défini le sujet de notre étude comme étant les communautés patrimoniales pirates : une organisation sociale existant par et pour une relation avec du patrimoine immatériel. Nous avons décidé d'étudier les expériences, les émotions et les actions des individus qui ont comme point commun l'utilisation et l'attachement à des services de bibliothèques et d'archives pirates. Nous avons également voulu étudier la conséquence du tournant numérique sur la question de la reproductibilité et de la propriété intellectuelle quand on la sort du contexte institutionnel et entrepreneurial. 

Nous nous sommes questionnés ainsi : Comment les dynamiques culturelles et les organisations sociales permettent-elles l'émergence d'une pratique archivistique populaire, autogestionnaire et clandestine ? Nous avons exploré la manière dont les individus découvrent et intègrent la culture et le fonctionnement des dépôts pirates. Nous avons ensuite cherché le point de rencontre entre l'expérience des individus et l'organisation de guérilla des bibliothèques pirates. Enfin, dans notre troisième chapitre, nous avons réalisé une typologie des différentes traductions des communautés patrimoniales pirates, du dépôt personnel aux bibliothèques d'Alexandrie clandestines contemporaines.\\

Une partie de la réponse se trouve au cœur de la nouvelle nature qu'ont acquise les objets culturels, avec la conséquence sur nos horizons mentaux de la massification de la capacité de copie des individus. Cette transformation de la capacité et du coût de copie entre en contradiction avec le mode de distribution et le fonctionnement économique de nos sociétés. Les pratiques de consommations culturelles actuelles sont perçues comme aliénantes et injustes. Où la culture et le divertissement, par extension sa patrimonialisation, ont été confisqués par un système marchand constitué du capitalisme centré autour de deux mannes de revenus principales : la surveillance et le service \footnote{Soshana Zuboff, \textit{L'âge du capitalisme de surveillance,}, éditions Zulma, 2019, \og Nouvel ordre économique qui revendique l'expérience humaine comme matière première gratuite à des fins de pratiques commerciales dissimulées d'extraction, de prédiction et de vente \fg }, les deux demandant un fonctionnement cloisonné, sécurisé. Cette mécanique est alors vécue comme une dépossession, qui pousse les individus à mettre en place un espace en dehors des codes sociaux classiques, un espace d'échange de l'information qui ne se limite qu'à la bande passante et à la capacité de stockage. Nous l'avons étudié chez nos enquêtés, tous trouvent un mécontentement, une injustice dans la manière dont le système culturel fonctionne, même s'ils n'y trouvent pas de considérations politiques, ce qui justifie leur passage à l'action en tant que pirates. Cette émotion, que l'on retrouve dans les écrits de D.Fabre, E. Bermès ou A. Kosnik. Ces dépôts pirates sont chargés d’autant de projets utopiques que d'espaces inquiétants et dangereux. Ces dépôts pirates sont chargés d’autant de projets utopiques que d'espaces inquiétants et dangereux. Le piratage offre un espace alternatif, qui, en plus de l’appât du gain économique, permet de toucher du doigt une forme de distribution alternative. Pour se retrouver dans les méandres des forums et des bibliothèques, il faut alors apprendre à maîtriser la culture, les interactions sociales, les nouveaux codes, les outils et l'esthétique particulière générée par la clandestinité et les représentations qu'ont les communautés d'elles-mêmes. Cet apprentissage est un phénomène social qui peut connecter des individus dans le monde physique (socialisation par les pairs, figure tutélaire), ou connecter l'individu à des ressources et des communautés qui permettent un apprentissage en autonomie. Une fois les règles et les outils maîtrisés, l'individu va mettre en place une routine et des habitudes de consommation incluant le piratage. Une rupture avec la pratique peut se produire si ces habitudes sont désorganisées (fermeture d'un site très utilisé, mise hors ligne d'un outil, moins de temps à dédier au piratage à cause d'un emploi, une situation économique qui permet de sortir du circuit de piratage). Une désorganisation des habitudes peut également pousser l'individu à évoluer et acquérir de nouvelles compétences. Ces compétences peuvent être également utiles et devenir une source de fierté dans d'autres interactions que la sphère culturelle pirate. L'individu peut alors devenir lui-même une passerelle pour ses pairs. \\

Deuxièmement, les bibliothèques pirates les plus connues comme Sci-Hub ou LibGen sont les éléments les plus visibles et médiatisés de ce que nous avons nommé une culture d'archivistique émergente. Cette culture tire son origine du lien émotionnel que les individus investissent dans le patrimoine culturel numérique. Ces dépôts sont à la fois le symbole de l'existence d'une fragilité concernant la conception de la propriété dans nos sociétés, et la tentative de créer une alternative, dans un espace non modéré. Cette dynamique de distribution culturelle parallèle a un impact extrêmement fort sur l'accès à la culture de toute une partie de la population la plus précaire, mais également des ramifications très profondes sur des pratiques professionnelles où des secteurs économiques entiers, comme par exemple le monde du développement de modèles d'intelligences artificielles génératives. En conclusion, l'étude du phénomène des bibliothèques pirates et des communautés patrimoniales pirates met en lumière les tensions inhérentes entre les modèles économiques actuels et la diffusion alternative de la culture, ainsi que les nouvelles compétences, solidarités et dangers qui en émergent.\\

De nombreux points traités dans ce travail pourraient être approfondis avec plus de temps. Nous n'avons pu nous entretenir qu'avec cinq enquêtés, et une grande partie de nos sources ne sont que quelques exemples perdus dans l'infinité du Web. Également, de nombreux projets d'exploitations n'ont pas été incorporés à notre travail. Nous avons également fait le choix de nous limiter au versant le plus proche de nos situations sociales, mais il existe toute une littérature sur les cultures pirates en dehors du nord économique. Nous n'avons pu également qu'effleurer certaines communautés comme les phénomènes low tech, les communautés DIY/Hacktivistes et les acteurs du \textit{Right to repair} et \textit{Right to Hack}. Ce travail ne peut pas proposer de réponses, seulement de nouvelles pistes à explorer. Pour prolonger les questionnements de ce mémoire après sa lecture, notre souhait est l'ouverture et l'approfondissement de la compréhension des pratiques auto-organisées et de leur influence sur l'accès à la culture, la conservation du patrimoine, mais également dans tous les domaines où les règles peuvent être tordues et réinventées pour des causes justes.  

\chapter{Sources}

\section*{Skyblogs}

Blog de Anti-telechargement, \href{https://web.archive.org/web/20130517194237/http://anti-telechargement.skyrock.com/}{https://web.archive.org/web/20130517194237/http://anti-telechargement.skyrock.com/}, créé en 2009, via Internet Archive. \\

Blog de Contre-piratages - Savoir Dire Non Aux Piratages!, \href{https://web.archive.org/web/20100428081928/http://contre-piratages.skyrock.com/}{https://web.archive.org/web/20100428081928/http://contre-piratages.skyrock.com/}, créé en 2009, via Internet Archive. \\

Blog de Film-a-Telecharger-Free, \href{https://web.archive.org/web/20130821121632/http://film-a-telecharger-free.skyrock.com/}{https://web.archive.org/web/20130821121632/http://film-a-telecharger-free.skyrock.com/}, créé en 2013,  via Internet Archive. \\

Blog de la-biblio-de-tetel - Entrez dans le monde de Tetel :3, \href{https://web.archive.org/web/20221202160656/https://la-biblio-de-tetel.skyrock.com/}{https://web.archive.org/web/20221202160656/https://la-biblio-de-tetel.skyrock.com/}, créé en 2012, via Internet Archive. \\

Blog de mangas-a-telecharger - mangas-a-telecharger, \href{https://web.archive.org/web/20130516071755/http://mangas-a-telecharger.skyrock.com/}{https://web.archive.org/web/20130516071755/http://mangas-a-telecharger.skyrock.com/}, créé en 2013, via Internet Archive. \\

Blog de Oo-Folie-des-Mangas-oO - Blog de Ecila-san, \href{https://web.archive.org/web/20110724003203/http://oo-folie-des-mangas-oo.skyrock.com/}{https://web.archive.org/web/20110724003203/http://oo-folie-des-mangas-oo.skyrock.com/}, créé en 2011, via Internet Archive. \\

Blog de telechargement09 - telechargement de music gratuit.system of a down.rammstein.Disturbed, \href{https://web.archive.org/web/20210504033421/https://telechargement09.skyrock.com/}{https://web.archive.org/web/20210504033421/https://telechargement09.skyrock.com/}, créé en 2007, via Internet Archive. \\

Blog de telechargerfilm - Film streaming vf, \href{https://web.archive.org/web/20180718101438/http://telechargerfilm.skyrock.com/}{https://web.archive.org/web/20180718101438/http://telechargerfilm.skyrock.com/}, créé en 2009, via Internet Archive. \\

Blog de Telechargergratuitement - ›› Téléchargement à volonté ! Logiciel gratuit à télécharger. Telecharger gratuitement les logiciels..., \href{https://web.archive.org/web/20181130030014/https://telechargergratuitement.skyrock.com/}{https://web.archive.org/web/20181130030014/https://telechargergratuitement.skyrock.com/}, créé en 2013, via Internet Archive. \\

Blog de telechargementgratos - Musiques à télécharger gratuitement, \href{https://web.archive.org/web/20130902023647/http://telechargementgratos.skyrock.com/}{https://web.archive.org/web/20130902023647/http://telechargementgratos.skyrock.com/}, créé en 2013, via Internet Archive. \\

Blog de le-monde-de-tetel - Bordel ambulant, \href{https://web.archive.org/web/20230711062347/https://le-monde-de-tetel.skyrock.com/}{https://web.archive.org/web/20230711062347/https://le-monde-de-tetel.skyrock.com/}, créé en 2013, via Internet Archive. \\

Blog de x-Yume-No-Manga-x - x-Yume-No-Manga-x Ou Le Paradis de L'Otaku', \href{https://web.archive.org/web/20230812201236/https://x-yume-no-manga-x.skyrock.com/}{https://web.archive.org/web/20230812201236/https://x-yume-no-manga-x.skyrock.com/}, créé en 2011, via Internet Archive.\\

\section*{Reddit}

\subsection*{r/Piracy}

u/ALIIERTx, \og By now it should be more moral to just pirate it \fg, 2024-06-10, \href{www.reddit.com/r/Piracy/comments/1dceyaz/by\_now\_it\_should\_be\_more\_moral\_to\_just\_pirate\_it/}{www.reddit.com/r/Piracy/comments/1dceyaz/by\_now\_it\_should\_be\_more\_moral\_to\_just\_pirate\_it/}. \\

u/JustHereForThePorn2x, \og When did you first become a pirate? And what was it \fg, 2024-02-20, \href{https://www.reddit.com/r/Piracy/comments/1avd288/when\_did\_you\_first\_become\_a\_pirate\_and\_what\_was\_it/}{https://www.reddit.com/r/Piracy/comments/1avd288/when\_did\_you\_first\_become\_a\_pirate\_and\_what\_was\_it/}. \\

\subsection*{r/DataHoarder}

u/ThePixelHunter, \og Court rules fan subtitles now illegal, piracy. Hoarding time? \fg, 2017-04-21, \href{www.reddit.com/r/DataHoarder/comments/66r52q/court\_rules\_fan\_subtitles\_now\_illegal\_piracy/}{www.reddit.com/r/DataHoarder/comments/66r52q/court\_rules\_fan\_subtitles\_now\_illegal\_piracy/}

u/Nath2125, \og Getting into archiving and data hoarding \fg, 2021-08-22, \href{www.reddit.com/r/DataHoarder/comments/p9a30g/getting\_into\_archiving\_and\_data\_hoarding/}{www.reddit.com/r/DataHoarder/comments/p9a30g/getting\_into\_archiving\_and\_data\_hoarding/} \\

u/Lee\_\_Jieun, \og Hoarding =/= Preservation \fg, 2022-10-11, \href{www.reddit.com/r/DataHoarder/comments/y1heob/hoarding\_preservation/}{www.reddit.com/r/DataHoarder/comments/y1heob/hoarding\_preservation/}. \\

u/pilimi\_anna, \og How to become a pirate archivist \fg, 2022-10-17, \href{www.reddit.com/r/DataHoarder/comments/y6pq40/how\_to\_become\_a\_pirate\_archivist/}{www.reddit.com/r/DataHoarder/comments/y6pq40/how\_to\_become\_a\_pirate\_archivist/}. \\ 

u/Titan\_91, \og I'm Archiving Bill Nye the Science Guy \fg, 2025-02-20, \href{www.reddit.com/r/DataHoarder/comments/1iubwd4/im\_archiving\_bill\_nye\_the\_science\_guy/}{www.reddit.com/r/DataHoarder/comments/1iubwd4/im\_archiving\_bill\_nye\_the\_science\_guy/}. \\

u/Invisibleflash, \og It takes \$\$ to be an archivist \fg, 2021-06-29, \href{www.reddit.com/r/DataHoarder/comments/oa8ir6/it\_takes\_to\_be\_an\_archivist/}{www.reddit.com/r/DataHoarder/comments/oa8ir6/it\_takes\_to\_be\_an\_archivist/}. \\

u/Mbox97, \og Manga Rock has officially shut down, I have an archive.\fg, 2020-01-04, \href{www.reddit.com/r/DataHoarder/comments/ek22j4/manga\_rock\_has\_officially\_shut\_down\_i\_have\_an/}{www.reddit.com/r/DataHoarder/comments/ek22j4/manga\_rock\_has\_officially\_shut\_down\_i\_have\_an/}. \\

u/Thynome, \og nHentai Archivist, a nhentai.net downloader suitable to save all of your favourite works before they're gone \fg, 2024-09-13, \href{www.reddit.com/r/DataHoarder/comments/1fg5yzy/nhentai\_archivist\_a\_nhentainet\_downloader/}{www.reddit.com/r/DataHoarder/comments/1fg5yzy/nhentai\_archivist\_a\_nhentainet\_downloader/}. \\

u/TheInternectivist, \og Our Role in the Preservation of the Cultural Heritage of Mankind \fg, 2017-05-21, \href{www.reddit.com/r/DataHoarder/comments/6cdsul/our\_role\_in\_the\_preservation\_of\_the\_cultural/}{www.reddit.com/r/DataHoarder/comments/6cdsul/our\_role\_in\_the\_preservation\_of\_the\_cultural/}. \\

u/jakuri69, \og PSA: Life is short. Don't spend too much time obsessively cataloguing your data collections. \fg, 2023-11-19, \href{www.reddit.com/r/DataHoarder/comments/17yylgg/psa\_life\_is\_short\_dont\_spend\_too\_much\_time/}{www.reddit.com/r/DataHoarder/comments/17yylgg/psa\_life\_is\_short\_dont\_spend\_too\_much\_time/}. \\

u/KyletheAngryAncap, \og So much will be lost. \fg, 2024-11-01, \href{www.reddit.com/r/DataHoarder/comments/1ghbc9g/so\_much\_will\_be\_lost/}{www.reddit.com/r/DataHoarder/comments/1ghbc9g/so\_much\_will\_be\_lost/}. \\

u/deleted, \og The Archivist's Dilemma \fg, 2019-08-04, \href{www.reddit.com/r/DataHoarder/comments/clz18v/the\_archivists\_dilemma/}{www.reddit.com/r/DataHoarder/comments/clz18v/the\_archivists\_dilemma/}. \\

u/AshleyUncia, \og The decline of 'Tech Literacy' having an influence on Data Hoarding. \fg, 2024-01-22, \href{www.reddit.com/r/DataHoarder/comments/19cx718/the\_decline\_of\_tech\_literacy\_having\_an\_influence/}{www.reddit.com/r/DataHoarder/comments/19cx718/the\_decline\_of\_tech\_literacy\_having\_an\_influence/}. \\

u/-Archivist, \og The Petabyte Porn Problem: Public webcam social interaction archiving. [Chaturbate, MyFreeCams, Streamate] \fg, 2017-04-13, \href{www.reddit.com/r/DataHoarder/comments/6583s2/the\_petabyte\_porn\_problem\_public\_webcam\_social/}{www.reddit.com/r/DataHoarder/comments/6583s2/the\_petabyte\_porn\_problem\_public\_webcam\_social/}

u/privacylawyer, \og The philosophy behind data hoarding and amateur archiving? \fg, 2018-05-01, \href{www.reddit.com/r/DataHoarder/comments/8gcji4/the\_philosophy\_behind\_data\_hoarding\_and\_amateur/}{www.reddit.com/r/DataHoarder/comments/8gcji4/the\_philosophy\_behind\_data\_hoarding\_and\_amateur/}. \\

u/-Archivist, \og The-Eye.eu: We now host the largest open repo of piracy metadata online today, with only more to come! \fg, 2019-11-24, \href{www.reddit.com/r/DataHoarder/comments/e0sb8b/theeyeeu\_we\_now\_host\_the\_largest\_open\_repo\_of/}{www.reddit.com/r/DataHoarder/comments/e0sb8b/theeyeeu\_we\_now\_host\_the\_largest\_open\_repo\_of/}.\\

u/Jakob4800, \og where is the line between data archival and piracy? \fg, 2022-09-21, \href{www.reddit.com/r/DataHoarder/comments/xjs9hm/where\_is\_the\_line\_between\_data\_archival\_and\_piracy/}{www.reddit.com/r/DataHoarder/comments/xjs9hm/where\_is\_the\_line\_between\_data\_archival\_and\_piracy/}. \\

u/-Archivist, \og Z-Library isn't really gone, but that maybe up to you. \fg, 2022-11-05, \href{www.reddit.com/r/DataHoarder/comments/ymiwzs/zlibrary\_isnt\_really\_gone\_but\_that\_maybe\_up\_to\_you/}{www.reddit.com/r/DataHoarder/comments/ymiwzs/zlibrary\_isnt\_really\_gone\_but\_that\_maybe\_up\_to\_you/}. \\

\section*{Twitter/X}

@fuyucchi.fr, \og Attention, ces quatre sites sont des sources très larges de livres et d'articles piratés, apprenez à les reconnaître pour mieux les éviter !!\fg, 2025-01-26. \href{https://bsky.app/profile/fuyucchi.fr/post/3lgo4xeycgc2c}{https://bsky.app/profile/fuyucchi.fr/post/3lgo4xeycgc2c}. \\

@desert\_\_varnish, \og They're both on my website now!! \fg, 2025-02-02, \href{https://x.com/desert\_\_varnish/status/1885833785721197009}{https://x.com/deser\_\_varnish/status/1885833785721197009}/ \\

@imitationlearn, \og can i just download (static?) websites (html+css?) and then keep them as a time capsule on my computer forever? \fg, 2025-03-04, \href{https://x.com/imitationlearn/status/1896887329182367937}{https://x.com/imitationlearn/status/1896887329182367937}. \\

@lunani2ta, \og new archivist starter pack just dropped \fg, 2025-03-04, \href{https://x.com/lunani2ta/status/1885183897480118655}{https://x.com/lunani2ta/status/1885183897480118655}. \\

@mattsutaaki, \og finally added a bio on soulseek \fg, 2025-02-13, \href{https://x.com/mattsutaaki/status/1890125617724785057}{https://x.com/mattsutaaki/status/1890125617724785057}. \\

@cycylcds, \og coucou j’ai téléchargé toutes les œuvres de Lynch (films et série) et je les ai mis dans un mega pour que ce soit plus facile pour toutes les personnes qui veulent (re)découvrir ce qu’il a fait \fg, 2025-01-18, \href{https://x.com/cycylcds/status/1880652858196132324}{https://x.com/cycylcds/status/1880652858196132324}. \\

\section*{Articles de presse}

Arthur Bayon, \og Enquête sur le «scantrad», le piratage de mangas à l'échelle industrielle \fg, le Figaro, 18 novembre 2024.  \\

Katy Guest, \og I can get any novel I want in 30 seconds': can book piracy be stopped ? \fg, The Guardian, 06 mars 2019. 

Le Figaro et l'AFP, \og Le Japon lance un programme, assisté par l’IA, pour lutter contre le piratage des mangas \fg, 08 décembre 2024. \\

Nicolas Gary, \og Longue vie aux "contenus culturels", qui alimentent le piratage \fg, Actualitté, le 04 septembre 2015. \\

Andrew Webster, \og An unprecedented Nntendo leak turns into a moral dilemma for archivists\fg, The Verge, 30 juillet 2020. \\




\chapter{Bibliographie}

\section*{Sentiments et communautés patrimoniales}

\subsection*{Monographies}
Bermès Emmanuelle, \guillemetleft De l'écran à l'émotion : Quand le numérique devient patrimoine \guillemetright, éditions de \textit{l'École nationale des chartes - PSL}, 242 pages, Paris, 2024 \\

De Kosnik Abigail, \guillemetleft Rogue Archives, Digital Cultural Memory and Media Fandom, éditions \textit{The MIT Press}, 430 pages, Cambridge, 2016 \\

Fabre Daniel, \og Émotions patrimoniales \fg, collection Ethnologie de la France,\textit{ éditions de la Maison des sciences de l’homme}, 2013. \\

Gray Jonathan, Sandvoss Cornel, Harrington C. Lee, \og Fandom: identities and communities in a mediated world \fg, éditions \textit{New York University Press}, 2017.\\

\subsection*{Articles de revues}
Blanchard Jean-François, \og Chiara Bortolotto, Le patrimoine culturel immatériel. Enjeux d’une nouvelle catégorie \fg, éditions \textit{ENS Editions}, 

Lessig Lawrence, \guillemetleft Industrie de la culture, pirates et \guillemetleft culture libre\guillemetright \guillemetright, dans \textit{Critique}, volume 73, n°6, p510-518, 2008\\

Paloque-Berges Camille, \og Vers des lieux de mémoire réticulaires ? \fg, dans \textit{Reset, Recherches en sciences sociales sur l'internet}, n°6, 2016. 

Paloque-Berges Camille et Schafer Valérie, \og Quand la communication devient patrimoine \fg, dans la Revue Hermès, \textit{éditions du CNRS}, p. 255-261, 2015. \\

\section*{Piratage et culture pirate}

\subsection*{ Monographies}

Baumgärtel Tilman, \guillemetleft A Reader on International Media Piracy : Pirate Essays \guillemetright, collection \textit{MediaMatters Series}, éditions \textit{Amsterdam University Press}, 2015\\

Paul Eve Martin, \og Warez, The Infrastructure and Aesthetics of Piracy \fg, éditions punctumbooks, 445 pages, 2021.\\

Scorpecci Danny et Stryszowski Piotr , \og Piracy of Digital Content \fg, éditions \textit{OECD Publishing}, 2009.

\subsection*{Articles de revues }

Bohannon John, \og Who's downloading pirated papers ? Everyone \fg, dans \textit{Science}, vol 352, p. 508-512, 2016. \\

Galluzzo Anthony, \og Longévité et résilience de l’accès illégal aux contenus culturels. Interroger les persistances et les sorties de carrières pirates \fg, dans \textit{Annales des Mines \- Gérer \& comprendre}, volume 145 n°3, p. 27-45, 2021. \\

Latrive Florent, \guillemetleft Du bon usage de la piraterie. Culture libre, science ouverte~\guillemetright, dans \textit{Multitude} volume 18, n° 4, p. 197-202, 2004 \\

Liang Lawrence, \guillemetleft Piratage, créativité et infrastructure : repenser l'accès à la culture \guillemetright, dans \textit{Tracés, revue de sciences humaines}, p183-202, 2014\\

Mattelart Tristan, \guillemetleft Piratage : apport et limites d'une infrastructure d'accès à la culture \guillemetright, dans \textit{Tracés, revue de sciences humaines}, n°26, p175-182, 2014\\

Tatchim Nicanor, \og Piratage, économie informelle et industries culturelles au Cameroun : approches socio-économique, communicationnelle et politique \fg, dans \textit{Études de communication, Médiations et usages sociaux des savoirs et de  l’information. Regards croisés France-Brésil}, n°57, 2021

Taylor Diana, \og Save As ... Knowledge and Transmission in the Age of Digital Technologies \fg, dans Imagining America n°7, 2010. 

Todd Darren, \og Pirate nation: how digital piracy is transforming business, society and culture \fg, éditions \textit{Kogan Page}, 2011. 

\section*{Archives et bibliothèques sur le web}

\subsection*{Monographies}

Brügger Niels, \og The archived web : doing history in the digital age \fg, éditions \textit{The MIT Press}, 2018. 

\subsection*{Articles de revues}

Magis Crisophe et Granjon Fabien, \og Numérique et libération de la produciton scientifique \fg, dans \textit{Variations. Revue internationale de théorie critique}, n°19, 2016. \\

Paul Eve Martin, \guillemetleft Lessons from the library : Extreme Minimalist Scaling at Pirate Ebook Platforms\guillemetright, dans \textit{Digital Humanities Quarterly}, dans la revue \textit{Providence}, volume 16, n° 2, 2022\\

\section*{Propriété intellectuelle et artistique}

\subsection*{Monographies}

National Research Council, \og The digital dilemma: intellectual property in the information age\og, \textit{National Academy Press}, 2010 \\

Sell Susan K,\og Private powern public law : the globalization of intellectual property rights \fg, collection Cambridge studies in international relations, éditions \textit{Cambridge University Press}, 2009

Goldstein Paul, \og, International Copyright : principles, law and practices \fg, \textit{Oxford University Press}, 2001. 
\chapter{Annexes}

\subsubsection*{Grille d'entretien \hrulefill}

        \subsection*{Question directrice : }

        \guillemetleft J'aimerais que vous me parliez de ce que vous piratez, depuis quand c'est un mode de consommation culturelle pour vous, que vous me racontiez comment vous le vivez au quotidien, si vous le partagez avec d'autres personnes et comment ça se traduit dans les différents domaines de la vie, professionnelle, personnelle, hobbies, dans vos études...\guillemetright

        \subsection*{Sujets que j'aimerais aborder : }
        
         \begin{itemize}

         \subsubsection*{Histoire personnelle :}
         
             \item Vous rappelez-vous de quand et de comment avez vous découvert le piratage ? Que recherchiez vous, que piratiez vous ?

            \item Saviez vous à l'époque que c'était du piratage ? Comment avez vous découvert que c'était du piratage ?

            \item Aujourd'hui, quelle est la fréquence à laquelle vous piratez du contenu culturel ?

            \item Vous rappelez-vous des périodes de votre vie où vous avez plus piraté ?
             
            \item Avez-vous aidé des proches moins à l'aise avec la technologie à accéder à des contenus piratés ? Avez-vous été dans une démarche de partage de cette pratique ?

        \subsubsection*{Questions sur les émotions :}
        
            \item Quelles sont les émotions que vous procure le fait d'avoir accès à un contenu culturel par le piratage ?

            \item Avez-vous déjà parlé de piratage à votre famille ? A vos ami.e.s ?

            \item Avez vous revendiqué une identité de pirate ou de hacker ?

            \item Qu'est ce qui vous a pousser à commencer à pirater du contenu culturel.
            
            \item Piratez vous pour les mêmes raisons qu'au début. 

        \subsubsection*{Questions techniques :}
        
             \item Avec quels outils piratiez vous (simple téléchargement, sites spécialisés, forums, discord, Reddit, torrenting, IPTV,  applications spécialisées)

             \item Est-ce que vous stockez les médias que vous piratez ? Comment stockez vous les médias que vous piratez ?

             \item Est ce que votre dépôt de contenu piraté suit une organisation particulière, ou un classement ?

             \item Est ce que votre dépôt de contenu piraté est accessible à distance, à vous ou à vos proches ?
             
         \end{itemize}

         pratique plus libre dans les années 2010 ?

\newpage

\subsubsection*{Pacte d'entretien d'information et de recueil de consentement}
    \hrulefill

    \subsection*{Question directrice}

    \guillemetleft J'aimerais que vous me parliez de ce que vous piratez, depuis quand c'est un mode de consommation culturelle pour vous, que vous me racontiez comment vous le vivez au quotidien, si vous le partagez avec d'autres personnes et comment ça se traduit dans les différents domaines de la vie, professionnelle, personnelle, dans vos hobbies, dans vos études...\guillemetright

    \subsection*{Collecte de données personnelles}

    Pour pouvoir contextualiser notre entretien, je vous interrogerai sur des sujets comme la catégorie socioprofessionnelle de vos parents ainsi que leur niveau d'études, votre catégorie socioprofessionnelle et votre parcours scolaire, votre utilisation des ordinateurs, des smartphones, et de votre niveau de maîtrise en informatique.\\

    Un enregistrement audio est également nécessaire à la production d'une transcription qui sera le matériel cité dans mon travail.

    \subsection*{Publication, modification et restitution}

    La transcription et les extraits utilisés seront publiés sous le nom que vous souhaitez, sous pseudonyme, ou avec une anonymisation totale.\\

    Vous pouvez également à tout moment rajouter, corriger, modifier ou supprimer un élément de l'entretien sur lequel vous souhaiteriez ajouter une précision, ou exprimer un changement d'avis.\\

    Comme ce travail se nourrit de vos expériences et de votre vécu, il est normal de vous laisser le choix d'avoir accès soit à une restitution des conclusions du travail final, soit au manuscrit du mémoire si cela vous intéresse.\\
    \vspace{0.5cm}
    
    Pour toutes questions ou modifications, vous pouvez me contacter à cette adresse :
    \vspace{0.1cm}
    
    jules.musquin@chartes.psl.eu
\subsubsection*{Message d'introduction à Ubuweb}

\begin{quote}
  \textit{  UbuWeb est une bibliothèque fantôme pirate composée de centaines de milliers d'artefacts d'avant-garde librement téléchargeables. Selon la lettre de la loi, le site est douteux ; nous violons ouvertement les normes du droit d'auteur et ne demandons presque jamais la permission. Presque tout ce qui se trouve sur le site est volé, déchiré et piqué à d'autres endroits, puis reposté. Nous n'avons jamais été poursuivis en justice - nous n'en sommes même pas proches. UbuWeb fonctionne sans argent - nous ne le prenons pas, nous ne le payons pas, nous n'y touchons pas ; vous ne trouverez jamais de publicité, de logo ou de boîte à dons. Nous n'avons jamais demandé de subvention ni accepté de parrainage ; nous restons heureusement non affiliés, ce qui nous permet de rester libres et propres, et de faire ce que nous voulons faire, de la manière dont nous voulons le faire. Plus important encore, UbuWeb a toujours été et sera toujours gratuit et ouvert à tous : il n'y a pas d'adhésion ou de mot de passe requis. Tout le travail est bénévole ; notre espace serveur et notre bande passante sont offerts par un groupe de gardiens intellectuels qui croient en l'accès libre à la connaissance. Économie du don et de l'abondance, avec un fort accent sur l'éducation globale, UbuWeb est visité quotidiennement par des dizaines de milliers de personnes de tous les continents. Nous figurons sur de nombreux programmes d'études, allant de ceux destinés aux enfants de maternelle qui étudient la poésie à motifs à ceux destinés aux étudiants de troisième cycle qui écoutent pendant des heures les séminaires de Jacques Lacan. Lorsque le site tombe en panne, comme c'est le cas pour la plupart des sites, nous sommes inondés de courriels de professeurs paniqués qui se demandent comment ils vont pouvoir donner leurs cours cette semaine-là.}

  \textit{About UbuWeb }; Excerpted \& adapted from \textit{Duchamp is My Lawyer: The Polemics, Pragmatics, and Poetics of UbuWeb}

  \end{quote}


\end{document}
